%%%%%%%%%%%%%%%%%%%%%%%%%%%%%%%%%%%%%%%%%%%%%%%%%%%%%%%%%%%%%%%%%%%%%%
% Lớp 12 | Nguyên hàm, tích phân | Nguyên hàm
% Dạng 2D4?1-2: Nguyên hàm cơ bản đa thức, phân thức
% Số câu: 112  (sinh tự động bởi tach_cau.py; từ nay thầy sửa trực tiếp file này)
%%%%%%%%%%%%%%%%%%%%%%%%%%%%%%%%%%%%%%%%%%%%%%%%%%%%%%%%%%%%%%%%%%%%%%

%%%%%%%%%%%%%%%% TRẮC NGHIỆM 4 PHƯƠNG ÁN %%%%%%%%%%%%%%%%

% nguồn: Phan1_Lop_12_Dai_so_Chuong_4
\begin{ex}%[2D4N1-2]%[ma:e0ab1cb2]
	Biết một nguyên hàm của hàm số $f(x)=\dfrac{1}{2x-1}$ là hàm  số $F(x)$ thỏa mãn $F(1)=5$. Khi đó $F(x)$ là hàm số nào sau đây?
	\choice 
		 {$F(x)=2\ln|2x-1|+5$}
		 {\True $F(x)=\dfrac{1}{2} \ln|2x-1|+5$}
		{$F(x)=\ln(2x-1)+5$} 
		{$F(x)=\ln|2x-1|+5$} 
	\loigiai{ 
		Ta có $F(x)=\displaystyle\int  \dfrac{1}{2x-1}\mathrm{~d} x=\dfrac{1}{2} \ln|2x-1|+C$ mà $F(1)=5\Rightarrow C=5$. \\
		Vậy $F(x)=\dfrac{1}{2} \ln|2x-1|+5$. 
	 }
\end{ex}

% nguồn: Phan1_Lop_12_Dai_so_Chuong_4
\begin{ex}%[2D4N1-2]%[ma:57a53469]
	Hàm  số nào sau đây là một nguyên hàm của hàm số $y=\dfrac{1}{x^2}$?
	\choice 
		 {$\dfrac{1}{x^3}$}
		 {\True $-\dfrac{1}{x}$}
		 {$\dfrac{1}{x}$}
		{$-\dfrac{1}{x^3}$} 
	\loigiai{ 
		Ta có $\left(-\dfrac{1}{x}\right)'=\dfrac{1}{x^2}$ 
	 }
\end{ex}

% nguồn: Phan1_Lop_12_Dai_so_Chuong_4
\begin{ex}%[2D4N1-2]%[ma:20f48716]
	Tìm nguyên hàm của  hàm số $f(x)=\dfrac{1}{5x-2}$.
	\choice 
		{\True  $\displaystyle\int \dfrac{d x}{5x-2}=\dfrac{1}{5} \ln|5x-2|+C$}
		{$\displaystyle\int  \dfrac{d x}{5x-2}=\ln|5x-2|+C$}
		{$\displaystyle\int  \dfrac{d x}{5x-2}=-\dfrac{1}{2} \ln|5x-2|+C$}
		{$\displaystyle\int \dfrac{d  x}{5x-2}=5\ln|5x-2|+C$}
	\loigiai{ 
		Ta có $\displaystyle\int \dfrac{d  x}{5x-2}=\dfrac{1}{5} \displaystyle\int \dfrac{\mathrm{~d}(5x-2)}{5x-2}=\dfrac{1}{5} \ln|5x-2|+C$.
	} 
\end{ex}

% nguồn: Phan1_Lop_12_Dai_so_Chuong_4
\begin{ex}%[2D4N1-2]%[ma:9ec93a60]
	Tìm nguyên hàm của  hàm số $f(x)=x^2+\dfrac{2}{x^2}$.
	\choice 
		{$\displaystyle\int f(x)\mathrm{~d} x=\dfrac{x^3}{3}+\dfrac{1}{x}+C$} 
		{\True $\displaystyle\int  f(x)\mathrm{~d} x=\dfrac{x^3}{3}-\dfrac{2}{x}+C$}
		 {$\displaystyle\int f(x)\mathrm{~d} x=\dfrac{x^3}{3}-\dfrac{1}{x}+C$}
		{$\displaystyle\int  f(x)\mathrm{~d} x=\dfrac{x^3}{3}+\dfrac{2}{x}+C$}
	\loigiai{ 
		Ta có $\displaystyle\int\left(x^2+\dfrac{2}{x^2}\right)\mathrm{~d} x=\dfrac{x^3}{3}-\dfrac{2}{x}+C$. 
	 }
\end{ex}

% nguồn: Phan1_Lop_12_Dai_so_Chuong_4
\begin{ex}%[2D4N1-2]%[ma:ab7d3a5f]
	 Họ nguyên hàm của hàm số $f(x)=\sin x+x$ là
	\choice 
		{\True $-\cos  x+\dfrac{1}{2}x^2+C$}
		{$\cos x+\dfrac{1}{2}x^2+C$} 
		 {$-\cos x-\dfrac{1}{2}x^2+C$}
		 {$\cos x+x^2+C$}
	 \loigiai{
		Ta có  $
		\displaystyle\int\left(\sin x+x\right)\mathrm{\,d}x=\displaystyle\int\sin x\,\mathrm{d}x+\displaystyle\int x\mathrm{\,d}x=-\cos x+\dfrac{1}{2}x^2+C$. 
	 }
\end{ex}

% nguồn: Phan1_Lop_12_Dai_so_Chuong_4
\begin{ex}%[2D4N1-2]%[ma:e60f933b]
	Cho hàm số $f(x)=x-\dfrac{1}{\sqrt{x}}$. Trong các hàm số dưới đây, hàm số nào  là một nguyên hàm của $f(x)$ trên $(0;+\infty)$?
	\choice 
		 {$F_1(x)=\dfrac{x^2}{2}+\sqrt{x}$}
		 {$F_2(x)=\dfrac{x^2}{2}-\sqrt{x}$}
		{$F_3(x)=\dfrac{x^2}{2}+2\sqrt{x}$} 
		{\True $F_4(x)=\dfrac{x^2}{2}-2\sqrt{x}$} 
	 \loigiai{
		 Ta có $F_1'(x)=x+\dfrac{1}{2\sqrt{x}}$; $F_2'(x)=x-\dfrac{1}{2\sqrt{x}}$; $ F_3'(x)=x+\dfrac{1}{\sqrt{x}}$; $F_4'(x)=x-\dfrac{1}{\sqrt{x}}$.\\
		Suy ra  $F_4'(x)=x-\dfrac{1}{\sqrt{x}}=f(x)$.\\
		Vậy $F_4(x)=\dfrac{x^2}{2}-2\sqrt{x}$ là một  nguyên hàm của $f(x)=x-\dfrac{1}{\sqrt{x}}$ trên $(0;+\infty)$.
	 }
\end{ex}

% nguồn: Phan1_Lop_12_Dai_so_Chuong_4
\begin{ex}%[2D4N1-2]%[ma:f445cfcd]
	 Nguyên hàm của hàm số $f(x)=x^3$ là
	\choice 
		{$x^4+C$} 
		{\True $\dfrac{x^4}{4}+C$} 
		 {$3x^2+C$}
		 {$\dfrac{x^2}{2}+C$}
	\loigiai{ 
		Ta có  $\displaystyle\int f(x)\mathrm{\,d}x=\displaystyle\int x^3 \mathrm{\,d}x=\dfrac{x^4}{4}+C$.
	 }
\end{ex}

% nguồn: Phan1_Lop_12_Dai_so_Chuong_4
\begin{ex}%[2D4N1-2]%[ma:82f349e8]
	Cho hàm số $f(x)=3x^2+2x$. Trong các hàm số dưới đây, hàm số nào là một  nguyên hàm của $f(x)$ trên $\mathbb{R}$?
	\choice 
		 {\True $F_1(x)=x^3+x^2-4$}
		 {$F_2(x)=\dfrac{x^3}{3}+\dfrac{x^2}{2}$}
		{$F_3(x)=x^3-x^2+1$} 
		 {$F_4(x)=3x^3+x^2$}
	\loigiai{ 
		 Ta có $F_1'(x)=3x^2+2x$; $F_2'(x)=x^2+x$; $F_3'(x)=3x^2-2x$; $F_4'(x)=9x^2+2x$.\\
		Suy ra  $F_1'(x)=3x^2+2x=f(x)$.\\
		Vậy $F_1(x)=x^3+x^2-4$ là một nguyên hàm của $f(x)=3x^2+2x$  trên $\mathbb{R}$.
	} 
\end{ex}

% nguồn: Phan1_Lop_12_Dai_so_Chuong_4
\begin{ex}%[2D4N1-2]%[ma:5c814988]
	Họ tất cả nguyên hàm của hàm số $f(x)=2x+4$  là
	\choice 
		 {$2x^2+4x+C$}
		{\True $x^2+4x+C$} 
		{$x^2+C$} 
		{$2x^2+C$} 
	\loigiai{ 
		Ta  có $\displaystyle\int f(x)\mathrm{~d} x=\displaystyle\int(2x+4)\mathrm{~d} x=x^2+4x+C$.
	} 
\end{ex}

% nguồn: Phan1_Lop_12_Dai_so_Chuong_4
\begin{ex}%[2D4N1-2]%[ma:84283d58]
	Tìm nguyên  hàm của hàm số $f(x)=x^{\sqrt{5}}$.
	\choice 
		{$\displaystyle\int f(x)\mathrm{~d} x=\dfrac{1}{\sqrt{5}-1} x^{\sqrt{5}-1}+C$} 
		{$\displaystyle\int f(x)\mathrm{~d}  x=x^{\sqrt{5}+1}+C$}
		{$\displaystyle\int  f(x)\mathrm{~d} x=\sqrt{5} x^{\sqrt{5}-1}+C$}
		{\True $\displaystyle\int  f(x)\mathrm{~d} x=\dfrac{1}{\sqrt{5}+1} x^{\sqrt{5}+1}+C$}
	 \loigiai{
		Áp dụng công  thức nguyên hàm cơ bản suy ra $\displaystyle\int f(x)\mathrm{~d} x=\dfrac{1}{\sqrt{5}+1} x^{\sqrt{5}+1}+C$.
	 }
\end{ex}

% nguồn: Phan1_Lop_12_Dai_so_Chuong_4
\begin{ex}%[2D4N1-2]%[ma:c67a7967]
	Họ tất cả các nguyên hàm của hàm  số $f(x)=x^2-2$ là
	\choice 
		{$\displaystyle\int  f(x)\mathrm{~d} x=2x+C$}
		 {\True $\displaystyle\int f(x)\mathrm{~d} x=\dfrac{x^3}{3}-2x+C$}
		{$\displaystyle\int  f(x)\mathrm{~d} x=x^2-2x+C$}
		{$\displaystyle\int f(x)\mathrm{~d} x=x^3-2x+C$} 
	\loigiai{ 
		Ta có $\displaystyle\int f(x)\mathrm{~d}  x=\displaystyle\int\left(x^2-2\right)\mathrm{~d} x=\displaystyle\int x^2\mathrm{~d} x-2\displaystyle\int \mathrm{~d} x=\dfrac{x^3}{3}-2x+C$.
	} 
\end{ex}

% nguồn: Phan1_Lop_12_Dai_so_Chuong_4
\begin{ex}%[2D4N1-2]%[ma:ab5d1b6e]
	Tìm nguyên hàm của hàm  số $f(x)=x^2+\dfrac{2}{x^2}$.
	 \choice
		 {$\dfrac{x^3}{3}-\dfrac{1}{x}+C$}
		{\True $\dfrac{x^3}{3}-\dfrac{2}{x}+C$} 
		{$\dfrac{x^3}{3}+\dfrac{1}{x}+C$} 
		{$\dfrac{x^3}{3}+\dfrac{2}{x}+C$} 
	\loigiai{ 
		Ta có $\displaystyle\int f(x)\mathrm{~d}  x=\displaystyle\int\left(x^2+\dfrac{2}{x^2}\right)\mathrm{~d} x=\dfrac{x^3}{3}-\dfrac{2}{x}+C$.
	 }
\end{ex}

% nguồn: Phan1_Lop_12_Dai_so_Chuong_4
\begin{ex}%[2D4N1-2]%[ma:5367f4e2]
	Cho $F(x)$ là một nguyên hàm của hàm  số $f(x)=\dfrac{1}{2x-3}$; biết $F(2)=1$. Giá trị $F(3)$ bằng
	 \choice
		{$F(3)=2\ln  3+1$}
		{\True $F(3)=\dfrac{1}{2}  \ln 3+1$}
		{$F(3)=\ln  3+1$}
		{$F(3)=\dfrac{1}{2}  \ln 3-1$}
	\loigiai{ 
		Ta có $F(x)=\displaystyle\int  f(x)\mathrm{~d} x=\displaystyle\int \dfrac{1}{2x-3} \mathrm{~d} x=\dfrac{1}{2} \ln|2x-3|+C$.\\
		Thay $x=2$ suy ra $C=1$ nên $F(x)=\dfrac{1}{2} \ln|2x-3|+1$. \\ 
		Vậy $F(3)=\dfrac{1}{2}  \ln 3+1$.
	 }
\end{ex}

% nguồn: Phan1_Lop_12_Dai_so_Chuong_4
\begin{ex}%[2D4N1-2]%[ma:6b816844]
	Cho hàm số $f(x)=3+\dfrac{1}{x}$. Trong các hàm số dưới đây, hàm số nào là một nguyên hàm của $f(x)$ trên  $(0;+\infty)$?
	 \choice
		{$F_1(x)=3x-\dfrac{1}{x^2}$} 
		{\True  $F_2(x)=3x+\ln x$}
		 {$F_3(x)=3x+\dfrac{1}{x^2}$}
		 {$F_4(x)=3x-\ln x$}
	 \loigiai{
		Ta có $F_1'(x)=3+\dfrac{2}{x^3}$; $F_2'(x)=3+\dfrac{1}{x}$; $F_3'(x)=3-\dfrac{2}{x^3}$;  $F_4'(x)=3-\dfrac{1}{x}$.\\
		 Suy ra $F_2'(x)=3+\dfrac{1}{x}=f(x)$.\\
		Vậy $F_2(x)=3x+\ln x$ là một nguyên hàm của $f(x)=3+\dfrac{1}{x}$ trên  $(0;+\infty)$.
	 }
\end{ex}

% nguồn: Phan1_Lop_12_Dai_so_Chuong_4
\begin{ex}%[2D4N1-2]%[ma:e719c1e0]
	Cho $a, b \in \mathbb{R},  a \neq 0$. Hàm số $y = \ln|ax+b|$ là một nguyên hàm của hàm số
	 \choice
		{$y =  \dfrac{1}{ax+b}$}
		{$y  = -\dfrac{1}{ax+b}$}
		{$y = \dfrac{-a}{ax+b}$} 
		 {\True $y = \dfrac{a}{ax+b}$}
	 \loigiai{
		Ta có $y' =  (\ln|ax+b|)'=\dfrac{a}{ax+b}$.
	} 
\end{ex}

% nguồn: Phan1_Lop_12_Dai_so_Chuong_4
\begin{ex}%[2D4N1-2]%[ma:217227b9]
	Tìm họ các nguyên hàm  $F(x)$ của hàm số $f(x)=\dfrac{x+3}{x+1}$.
	 \choice
		{$F(x)=x-\ln|x+1|+C$} 
		{$F(x)=x+\ln|x+1|+C$} 
		 {$F(x)=x-3\ln|x+1|+C$}
		{\True $F(x)=x+2\ln|x+1|+C$} 
	\loigiai{Ta có $\displaystyle\int f(x)\mathrm{~d} x=\displaystyle\int \dfrac{x+3}{x+1} \mathrm{~d}  x=\displaystyle\int \dfrac{x+1+2}{x+1} \mathrm{~d} x=\displaystyle\int\left(1+\dfrac{2}{x+1}\right)\mathrm{~d} x=x+2\ln|x+1|+C$.}
\end{ex}

% nguồn: Phan1_Lop_12_Dai_so_Chuong_4
\begin{ex}%[2D4N1-2]%[ma:c34d705f]
	 Cho hàm số $f(x) = e^{2x}$. Khi đó
	\choice 
		{$\displaystyle\int f(x)  \mathrm{d}x = e^{2x} + C$}
		{$\displaystyle\int  f(x) \mathrm{d}x = 2e^{2x} + C$}
		 {\True $\displaystyle\int f(x) \mathrm{d}x = \dfrac{1}{2}e^{2x} + C$}
		{$\displaystyle\int f(x) \mathrm{d}x = 2e^x + C$} 
	 \loigiai{
		Ta có $\displaystyle\int f(x) \mathrm{d}x = \displaystyle\int e^{2x} \mathrm{d}x=\displaystyle\int \left( e^{2}\right)^x \mathrm{d}x= \dfrac{e^{2x}}{\ln  e^{2}} + C= \dfrac{1}{2}e^{2x} + C$.\\
		Vậy phương án đúng là  $\displaystyle\int f(x) \mathrm{d}x = \dfrac{1}{2}e^{2x} + C$.
	} 
\end{ex}

% nguồn: Phan1_Lop_12_Dai_so_Chuong_4
\begin{ex}%[2D4N1-2]%[ma:3339c366]
	Biết $\displaystyle\int\left(\dfrac{1}{2x}+x^5\right)\mathrm{~d} x=a \ln|x|+b x^6+C$ với $(a, b \in Q, C\in \mathbb{R})$. Tính  $a^2+b$.
	 \choice
		{\True  $\dfrac{5}{12}$}
		{$9$} 
		 {$\dfrac{7}{13}$}
		 {$\dfrac{7}{6}$}
	 \loigiai{
		Ta có $\displaystyle\int\left(\dfrac{1}{2x}+x^5\right)\mathrm{~d} x=\dfrac{1}{2}  \ln|x|+\dfrac{1}{6} x^6+C$ nên $a=\dfrac{1}{2}$, $b=\dfrac{1}{6} \Rightarrow a^2+b=\dfrac{1}{4}+\dfrac{1}{6}=\dfrac{5}{12}$.
	 }
\end{ex}

% nguồn: Phan1_Lop_12_Dai_so_Chuong_4
\begin{ex}%[2D4N1-2]%[ma:bd597830]
	Cho hàm số $y=F(x)$ là một nguyên hàm  của hàm số $y=x^2$. Tính $F'(25)$.
	\choice 
		 {$25$}
		 {$125$}
		{$5$} 
		{\True $625$} 
	\loigiai{ 
		Ta có $F(x)$ là  một nguyên hàm của hàm số $y=x^2$ nên ta có $F'(x)=x^2, \forall x \in \mathbb{R}$.\\
		Do đó $F'(25)=25^2=625$. 
	 }
\end{ex}

% nguồn: Phan1_Lop_12_Dai_so_Chuong_4
\begin{ex}%[2D4N1-2]%[ma:8a61b9a4]
	Họ tất cả các nguyên hàm của hàm số $f(x)=5x^4-8x^3-6x$ là 
	 \choice
		{\True  $F(x)=x^5-2x^4-3x^2+C$}
		{$F(x)=x^5-x^4-x^2+C$} 
		{$F(x)=x^5-4x^4-2x^2+C$} 
		{$F(x)=x^5+2x^4-3x^2+C$} 
	 \loigiai{
		Ta có $F(x)=\displaystyle\int\left(5x^4-8x^3-6x\right)\mathrm{~d} x=x^5-2x^4-3x^2+C$.\\ 
		Vậy một nguyên hàm  của hàm số $f(x)$ là $F(x)=x^5-2x^4-3x^2+C$.
	} 
\end{ex}

% nguồn: Phan1_Lop_12_Dai_so_Chuong_4
\begin{ex}%[2D4N1-2]%[ma:0166b5f7]
	Tìm họ nguyên hàm  của hàm số $f(x)=\dfrac{x^2-x+1}{x-1}$.
	 \choice
		{$x+\dfrac{1}{x-1}+C$} 
		 {$1+\dfrac{1}{(x-1)^2}+C$}
		{\True  $\dfrac{x^2}{2}+\ln|x-1|+C$}
		{$x^2+\ln|x-1|+C$} 
	\loigiai{ 
		Ta có \begin{eqnarray*}
		\displaystyle\int \dfrac{x^2-x+1}{x-1}  \mathrm{~d} x& =&\displaystyle\int \dfrac{\left(x^2-x\right)+1}{x-1} \mathrm{~d} x=\displaystyle\int x \mathrm{~d} x+\displaystyle\int \dfrac{1}{x-1} \mathrm{~d} x\\
		&=&  \displaystyle\int x \mathrm{~d} x+\displaystyle\int \dfrac{1}{x-1} \mathrm{~d}(x-1)=\dfrac{1}{2} x^2+\ln|x-1|+C.
	\end{eqnarray*}
} 
\end{ex}

% nguồn: Phan1_Lop_12_Dai_so_Chuong_4
\begin{ex}%[2D4N1-2]%[ma:32a27657]
	Giá trị  của $\displaystyle\int_{-2}^{1}(4x^3+3x^2+8x)\mathrm{\,d}x+\displaystyle\int_{1}^{2}(4x^3+3x^2+8x)\mathrm{\,d}x$ bằng
	 \choice
		{\True $16$} 
		 {$-16$}
		 {$52$}
		 {$0$}
	 \loigiai{
		 $\displaystyle\int_{-2}^{1}(4x^3+3x^2+8x)\mathrm{\,d}x+\displaystyle\int_{1}^{2}(4x^3+3x^2+8x)\mathrm{\,d}x=\displaystyle\int_{-2}^{2}(4x^3+3x^2+8x)\mathrm{\,d}x=(x^4+x^3+4x^2)\bigg|_{-2}^{2}=40-24=16$
	} 
\end{ex}

% nguồn: Phan1_Lop_12_Dai_so_Chuong_4
\begin{ex}%[2D4N1-2]%[ma:caedd75e]
	Họ các nguyên hàm của  hàm số $y=x(x+1)^5$ là
	 \choice
		{$\dfrac{(x+1)^7}{7}+\dfrac{(x+1)^6}{6}+C$} 
		 {$6(x+1)^5+5(x+1)^4+C$}
		{$6(x+1)^5-5(x+1)^4+C$} 
		{\True $\dfrac{(x+1)^7}{7}-\dfrac{(x+1)^6}{6}+C$} 
	\loigiai{ 
		 Ta có $y=x(x+1)^5=(x+1-1)(x+1)^5=(x+1)^6-(x+1)^5$.\\
		 Kiểm tra $F(x)=\dfrac{(x+1)^7}{7}-\dfrac{(x+1)^6}{6}+C$.\\
		 Ta có $F'(x)=\dfrac{7(x+1)^6}{7}-\dfrac{6(x+1)^5}{6}=(x+1)^6-(x+1)^5$.
	} 
\end{ex}

% nguồn: Phan1_Lop_12_Dai_so_Chuong_4
\begin{ex}%[2D4N1-2]%[ma:db7eb748]
	Họ nguyên hàm của hàm  số $f(x)=4x^3-2\,024$ là
	 \choice
		{\True  $x^4-2\,024x+C$}
		 {$4x^3-2\,024x+C$}
		{$12x^3+C$} 
		{$x^4+C$} 
	\loigiai{ 
		Ta có $\displaystyle\int\left(4x^3-2\,024\right)\mathrm{~d}  x=x^4-2\,024x+C$.
	} 
\end{ex}

% nguồn: Phan1_Lop_12_Dai_so_Chuong_4
\begin{ex}%[2D4N1-2]%[ma:84b8a7c0]
	Khẳng định nào sau đây {\bf{sai}}? 
	\choice 
		{$\displaystyle\int 0\mathrm{\,d}x=C$} 
		{$\displaystyle\int x^4\mathrm{\,d}x=\dfrac{x^5}{5}+C$} 
		 {\True $\displaystyle\int \dfrac{1}{x} \mathrm{\,d}x=\ln x+C$}
		{$\displaystyle\int \mathrm{e}^x \mathrm{\,d}x=\mathrm{e}^x+C$} 
	\loigiai{ 
		Ta có  $\displaystyle\int \dfrac{1}{x} \mathrm{\,d}x=\ln|x|+C$.
	 }
\end{ex}

% nguồn: Phan1_Lop_12_Dai_so_Chuong_4
\begin{ex}%[2D4N1-2]%[ma:d358ebb1]
	Họ nguyên hàm của hàm số $f(x)=e^x$  là
	\choice 
		{$x e^{x-1}+C$} 
		{\True  $e^x+C$}
		{$\dfrac{1}{e^x}+C$} 
		{$\dfrac{e^{x+1}}{x+1}+C$} 
	\loigiai{ 
		Ta có $\displaystyle \int e^x  \mathrm{\,d}x=e^x+C $.
	} 
\end{ex}

% nguồn: Phan1_Lop_12_Dai_so_Chuong_4
\begin{ex}%[2D4N1-2]%[ma:eeaeb0e3]
	Tìm họ  nguyên hàm của hàm số $ y=x^{2025} $.
	 \choice
		 {$2026x^{2025}+C$}
		 {$x^{2026}+C$}
		{$2026x^{2026}+C$} 
		{\True  $\dfrac{1}{2026}x^{2026}+C$}
	\loigiai{ 
		Ta có $\displaystyle\int  x^{2025} \mathrm{\,d}x=\dfrac{1}{2026}x^{2026}+C$.
	 }
\end{ex}

% nguồn: Phan1_Lop_12_Dai_so_Chuong_4
\begin{ex}%[2D4N1-2]%[ma:334b5858]
	Khẳng định  nào đúng?
	 \choice
		{\True  $\displaystyle\int\left(x-\dfrac{1}{x}\right)^2\mathrm{\,d}x=\dfrac{x^3}{3}-2x-\dfrac{1}{x}+C$}
		{$\displaystyle\int\left(x-\dfrac{1}{x}\right)^2\mathrm{\,d}x=\dfrac{x^3}{3}-2x+\dfrac{1}{x}+C$} 
		{$\displaystyle\int\left(x-\dfrac{1}{x}\right)^2\mathrm{\,d}x=\dfrac{1}{3}\left(x-\dfrac{1}{x}\right)^3+C$} 
		 {$\displaystyle\int\left(x-\dfrac{1}{x}\right)^2\mathrm{\,d}x=\dfrac{1}{3}\left(x-\dfrac{1}{x}\right)^3\left(1+\dfrac{1}{x^2}\right)+C$}
	 \loigiai{
		Ta  có $\displaystyle\int\left(x-\dfrac{1}{x}\right)^2\mathrm{\,d}x=\displaystyle\int\left(x^2-2+\dfrac{1}{x^2}\right)\mathrm{\,d}x=\dfrac{x^3}{3}+2x-\dfrac{1}{x}+C$
	 }
\end{ex}

% nguồn: Phan1_Lop_12_Dai_so_Chuong_4
\begin{ex}%[2D4N1-2]%[ma:c0d08488]
	Cho số thực $\alpha \ne -1$. Hàm số nào sau  đây là một nguyên hàm của hàm số $y = x^\alpha$?
	 \choice
		{$y =  x^{\alpha+1}$}
		{$y = (\alpha+1)x^{\alpha+1}$} 
		{$y  = \dfrac{x^\alpha}{\alpha+1}$}
		{\True  $y = \dfrac{x^{\alpha+1}}{\alpha+1}$}
	\loigiai{ 
		Ta có $\displaystyle\int x^\alpha\mathrm{\,d}x=\dfrac{x^{\alpha+1}}{\alpha+1}+C$.\\ 
		Vậy  $y=\dfrac{x^{\alpha+1}}{\alpha+1}$.
	 }
\end{ex}

% nguồn: Phan1_Lop_12_Dai_so_Chuong_4
\begin{ex}%[2D4N1-2]%[ma:4f137f4a]
	Cho $F(x)$ là một nguyên hàm của $f(x)=\dfrac{2}{x+2}$. Biết $F(-1)=1$, khi đó $F(2)$  bằng
	\choice 
		{$2\ln  3+2$}
		 {\True $4\ln 2+1$}
		{$\ln 8+1$} 
		 {$2\ln 4$}
	 \loigiai{
		Ta có  $F(x)=\displaystyle\int f(x)\mathrm{~d} x=\displaystyle\int \dfrac{2}{x+2} \mathrm{~d} x=2\ln|x+2|+C$.\\
		Do $F(-1)=1$  nên $2\ln|-1+2|+C=1\Leftrightarrow C=1$.\\
		Vậy $F(x)=2\ln|x+2|+1\Rightarrow F(2)=2\ln 4+1=4\ln 2+1$. 
	} 
\end{ex}

% nguồn: Phan1_Lop_12_Dai_so_Chuong_4
\begin{ex}%[2D4N1-2]%[ma:8e583fa7]
	Hàm số nào sau đây là một  nguyên hàm của hàm số $y=x^4$?
	\choice 
		{$-\dfrac{x^5}{5}$} 
		{$4x^3$} 
		 {\True$\dfrac{x^5}{5}+1$}
		 {$-4x^3-1$}
	 \loigiai{
		 Ta có $\left(\dfrac{x^5}{5}+1\right)'=x^4$.
	 }
\end{ex}

% nguồn: Phan1_Lop_12_Dai_so_Chuong_4
\begin{ex}%[2D4N1-2]%[ma:39b4f87b]
	Họ  nguyên hàm của hàm số $f(x)=3x^2+2x+5$ là
	\choice 
		{\True  $x^3+x^2+5x+C$}
		{$x^3+x+C$} 
		{$x^3+x^2+C$} 
		 {$x^3+x^2+5$}
	 \loigiai{
		$\displaystyle\int f(x)\mathrm{\,d}x=\displaystyle\int(3x^2+2x+5)\mathrm{\,d}x=x^3+x^2+5x+C$ 
	 }
\end{ex}

% nguồn: Phan1_Lop_12_Dai_so_Chuong_4
\begin{ex}%[2D4N1-2]%[ma:15c81791]
	Tìm nguyên hàm của hàm số $\displaystyle\int\left(x^2+\dfrac{3}{x}-2\sqrt{x}\right)\mathrm{~d}  x$.
	 \choice
		{$\dfrac{x^3}{3}+3\ln|x|+\dfrac{4}{3}  \sqrt{x^3}+C$}
		{\True $\dfrac{x^3}{3}+3\ln|x|-\dfrac{4}{3}  \sqrt{x^3}+C$}
		{$\dfrac{x^3}{3}+3\ln x-\dfrac{4}{3}  \sqrt{x^3}$}
		{$\dfrac{x^3}{3}-3\ln|x|-\dfrac{4}{3} \sqrt{x^3}+C$} 
	\loigiai{ 
		 Ta có $\displaystyle\int\left(x^2+\dfrac{3}{x}-2\sqrt{x}\right)\mathrm{~d} x=\dfrac{1}{3} x^3+3\ln|x|-\dfrac{4}{3} \sqrt{x^3}+C$.
	} 
\end{ex}

% nguồn: Phan1_Lop_12_Dai_so_Chuong_4
\begin{ex}%[2D4N1-2]%[ma:5861b294]
	Họ tất cả các nguyên hàm của hàm số $f(x)=2x+4$ là 
	 \choice
		 {$2x^2+4x+C$}
		 {\True $x^2+4x+C$}
		{$x^2+C$} 
		 {$2x^2+C$}
	 \loigiai{
		Ta  có $\displaystyle\int f(x)\mathrm{~d} x=\displaystyle\int(2x+4)\mathrm{~d} x=2\displaystyle\int x \mathrm{~d} x+4\displaystyle\int \mathrm{~d} x=x^2+4x+C$.
	} 
\end{ex}

% nguồn: Phan1_Lop_12_Dai_so_Chuong_4
\begin{ex}%[2D4N1-2]%[ma:4d54d56d]
	Họ nguyên  hàm của hàm số $f(x)=x-\dfrac{1}{x}$ là
	\choice 
		{\True  $\dfrac{1}{2}x^2-\ln\left|x\right|+C$}
		 {$\dfrac{1}{2}x^2-\ln\left|x\right|$}
		 {$\dfrac{1}{2}x^2-\ln x+C$}
		 {$1-\ln\left|x\right|+C$}
	\loigiai{ 
		 Ta có $\displaystyle\int\left(x-\dfrac{1}{x}\right)\mathrm{\,d}x=\displaystyle\int x\,\mathrm{d}x-\displaystyle\int\dfrac{1}{x}\mathrm{\,d}x=\dfrac{1}{2}x^2-\ln\left|x\right|+C$.
	} 
\end{ex}

% nguồn: Phan1_Lop_12_Dai_so_Chuong_4
\begin{ex}%[2D4N1-2]%[ma:755ce36f]
	Tìm nguyên hàm của hàm số $f(x)=\sqrt{2x-1}$. 
	 \choice
		 {\True $\displaystyle\int f(x)\mathrm{~d} x=\dfrac{1}{3}(2x-1) \sqrt{2x-1}+C$}
		{$\displaystyle\int f(x)\mathrm{~d} x=\dfrac{1}{2}  \sqrt{2x-1}+C$}
		 {$\displaystyle\int f(x)\mathrm{~d} x=\dfrac{2}{3}(2x-1) \sqrt{2x-1}+C$}
		{$\displaystyle\int f(x)\mathrm{~d} x=-\dfrac{1}{3} \sqrt{2x-1}+C$} 
	\loigiai{ 
		\begin{eqnarray*}
			\displaystyle\int f(x)\mathrm{~d} x& =&\displaystyle\int \sqrt{2x-1} \mathrm{~d} x \\ 
			&=&\displaystyle\int  \dfrac{(2x-1)^{\tfrac{1}{2}} \mathrm{~d}(2x-1)}{2}\\
			&=&\dfrac{1}{2} \cdot \dfrac{(2x-1)^{\tfrac{3}{2}}}{\dfrac{3}{2}}+C\\ 
			 &=&\dfrac{1}{3}(2x-1) \sqrt{2x-1}+C.
		\end{eqnarray*}
	 }
\end{ex}

% nguồn: Phan1_Lop_12_Dai_so_Chuong_4
\begin{ex}%[2D4N1-2]%[ma:36ae3331]
	Khẳng định nào  sau đây {\bf sai}?
	\choice 
		{$\displaystyle\int\left[f(x)+x\right]  \mathrm{\,d}x=\displaystyle\int f(x) \mathrm{\,d}x+\displaystyle\int x \mathrm{\,d} x$}
		{$\displaystyle\int\left[f(x)-1\right] \mathrm{\,d}x=\displaystyle\int f(x) \mathrm{\,d}x-\displaystyle\int  1\mathrm{\,d} x$}
		{\True $\displaystyle\int k f(x) \mathrm{\,d}x=k \displaystyle\int f(x) \mathrm{\,d}x,  \forall k \in \mathbb{R}$}
		{$\displaystyle\int k f(x) \mathrm{\,d}x=k \displaystyle\int f(x) \mathrm{\,d}x, \forall k \in \mathbb{R}, k  \neq 0$}
	 \loigiai{
		$\displaystyle\int k f(x) \mathrm{\,d}x=k \displaystyle\int f(x) \mathrm{\,d}x, \forall k \in \mathbb{R}$ là khẳng định sai vì $k\neq 0$. 
	 }
\end{ex}

% nguồn: Phan1_Lop_12_Dai_so_Chuong_4
\begin{ex}%[2D4N1-2]%[ma:d6341791]
	 Họ nguyên hàm của hàm số $f(x)=\cos x-1$ là
	\choice 
		{\True  $\sin x-x+C$}
		 {$\sin x+C$}
		{$-\sin x+C$} 
		{$-\sin x-x+C$} 
	 \loigiai{
		Ta có $\displaystyle\int f(x)\mathrm{\,d}x=\int(\cos  x-1)\mathrm{\,d}x=\sin x-x+C$.
	 }
\end{ex}

% nguồn: Phan1_Lop_12_Dai_so_Chuong_4
\begin{ex}%[2D4N1-2]%[ma:3e984858]
	Khẳng định nào dưới đây đúng? 
	\choice 
		 {$\displaystyle \int\left(4 x^3-2\right)\mathrm{\,d}x=x^4+C$}
		 {$\displaystyle \int\left(4 x^3-2\right)\mathrm{\,d}x=12 x^2+C$}
		{\True  $\displaystyle \int\left(4 x^3-2\right)\mathrm{\,d}x=x^4-2 x+C$}
		{$\displaystyle \int\left(4  x^3-2\right)\mathrm{\,d}x=12 x^4-2 x+C$}
	 \loigiai{
		 Ta có $\displaystyle \int\left(4 x^3-2\right)\mathrm{\,d}x=4 \cdot \dfrac{x^4}{4}-2 x+C=x^4-2 x+C$.
	 }
\end{ex}

% nguồn: Phan1_Lop_12_Dai_so_Chuong_4
\begin{ex}%[2D4N1-2]%[ma:4d0eadb0]
	Nguyên hàm $F(x)$ của hàm  số $f(x)=\dfrac{1}{3} x^3-2x^2+x-2024$ thỏa mãn $F(1)=-2024$ là
	\choice 
		 {$\dfrac{1}{12} x^4-\dfrac{2}{3} x^3+\dfrac{1}{2} x^2-2024x+\dfrac{5}{12}$}
		{$\dfrac{1}{12} x^4-\dfrac{2}{3} x^3+\dfrac{1}{2}  x^2-2024x-\dfrac{1}{12}$}
		{\True $\dfrac{1}{12} x^4-\dfrac{2}{3} x^3+\dfrac{1}{2} x^2-2024x+\dfrac{1}{12}$} 
		 {$\dfrac{1}{12} x^4-\dfrac{2}{3} x^3+\dfrac{1}{2} x^2-2024x-\dfrac{5}{12}$}
	 \loigiai{
		Ta  có
		\begin{eqnarray*}
			\displaystyle\int\left(\dfrac{1}{3} x^3-2x^2+x-2024\right)\mathrm{~d} x& =& \dfrac{1}{3} \displaystyle\int x^3\mathrm{~d} x-2\displaystyle\int x^2\mathrm{~d} x+\displaystyle\int  x \mathrm{~d} x-2024\displaystyle\int \mathrm{~d} x\\
			 &=&\dfrac{1}{3} \cdot \dfrac{x^4}{4}-2\cdot \dfrac{x^3}{3}+\dfrac{x^2}{2}-2024x+C\\
			 &=&\dfrac{1}{12} x^4-\dfrac{2}{3} x^3+\dfrac{1}{2} x^2-2024x+C.
		\end{eqnarray*}
		$F(1)=-2024\Leftrightarrow \dfrac{1}{12}  \cdot 1^4-\dfrac{2}{3} \cdot 1^3+\dfrac{1}{2} \cdot 1^2-2024\cdot 1+C=-2024\Leftrightarrow C=\dfrac{1}{12}$.\\
	Vậy nguyên hàm $F(x)$  của hàm số $f(x)=\dfrac{1}{3} x^3-2x^2+x-2024$ thỏa mãn $F(1)=-2024$ là $\dfrac{1}{12} x^4-\dfrac{2}{3} x^3+\dfrac{1}{2} x^2-2024x+\dfrac{1}{12}$.}
\end{ex}

% nguồn: Phan1_Lop_12_Dai_so_Chuong_4
\begin{ex}%[2D4N1-2]%[ma:4a9abadb]
	Một nguyên hàm của hàm  số $f(x)=x^3+2x$ có dạng $F(x)=a x^4+b x^2$. Tính $T=4a+b$.
	\choice 
		{$T=0$} 
		 {$T=1$}
		{\True $T=2$} 
		 {$T=3$}
	\loigiai{ 
		Ta có $F(x)=a x^4+b x^2$ là nguyên hàm của  hàm số $f(x)=x^3+2x$ nên
		$F'(x)=f(x),  \forall x \in \mathbb{R}$\\
		Khi đó $4a x^3+2b  x=x^3+2x, \forall x \in \mathbb{R}\Leftrightarrow\heva{&4a=1\\
		&2b=2}  \Leftrightarrow\heva{&a=\dfrac{1}{4} \\
		 &b=1}$\\
		Do đó $T=4a+b=4\cdot \dfrac{1}{4}+1=2$. 
	 }
\end{ex}

% nguồn: Phan1_Lop_12_Dai_so_Chuong_4
\begin{ex}%[2D4N1-2]%[ma:a3bd12b3]
	 Nguyên hàm $F(x)$ của hàm số $f(x)=(2x-3)^2$ thỏa mãn $F(-1)=-17$ là
	 \choice
		{$\dfrac{(2x-3)^3}{3}+\dfrac{4}{3}$} 
		{\True $\dfrac{4}{3} x^3-6x^2+9x-\dfrac{2}{3}$} 
		{$\dfrac{4}{3} x^3-6x^2+9x+\dfrac{8}{3}$} 
		 {$\dfrac{4}{3} x^3-6x^2+9x+\dfrac{2}{3}$}
	 \loigiai{
		Ta có 
		\begin{eqnarray*}
			\displaystyle\int(2x-3)^2\mathrm{~d} x& =&  \displaystyle\int\left(4x^2-12x+9\right)\mathrm{~d} x=\\
			&=&4\displaystyle\int x^2\mathrm{~d} x-12\displaystyle\int x \mathrm{~d}  x+9\displaystyle\int \mathrm{~d} x\\
			 &=&\dfrac{4}{3} x^3-6x^2+9x+C.
		\end{eqnarray*}
		$F(-1)=-17\Leftrightarrow \dfrac{4}{3}  \cdot(-1)^3-6\cdot(-1)^2+9\cdot(-1)+C=-17\Leftrightarrow C=-\dfrac{2}{3}$.\\
		Vậy nguyên hàm $F(x)$ của hàm số $f(x)=(2x-3)^2$ thỏa mãn $F(-1)=-17$ là $\dfrac{4}{3} x^3-6x^2+9x-\dfrac{2}{3}$. 
	 }
\end{ex}

% nguồn: Phan1_Lop_12_Dai_so_Chuong_4
\begin{ex}%[2D4N1-2]%[ma:e7bb21b7]
	Trong các hàm số sau, hàm số nào có một nguyên hàm là  hàm số $F(x)=\ln |x|$?
	 \choice
		{$f(x)=x$} 
		 {\True $f(x)=\dfrac{1}{x}$}
		 {$f(x)=\dfrac{x^3}{3}$}
		 {$f(x)=|x|$}
	\loigiai{ 
		Áp dụng công thức  SGK.
	 }
\end{ex}

% nguồn: Phan1_Lop_12_Dai_so_Chuong_4
\begin{ex}%[2D4N1-2]%[ma:6f267cb4]
	$\displaystyle\int\left(x^2+1\right)\mathrm{\,d}x$  bằng
	\choice 
		{$2x+C$} 
		{\True $\dfrac{1}{3}x^3+x+C$} 
		{$x^3+x+C$} 
		{$3x^3+x+C$} 
	 \loigiai{
		Ta  có
	\[\displaystyle\int\left(x^2+1\right)\mathrm{\,d}x=\dfrac{1}{3}x^3+x+C.\]} 
\end{ex}

% nguồn: Phan1_Lop_12_Dai_so_Chuong_4
\begin{ex}%[2D4N1-2]%[ma:49bdf450]
	Hàm số nào sau đây là một nguyên  hàm của hàm số $ y=x^5$?
	 \choice
		 {$ y=x^6$}
		{$ y=5x^4$} 
		{\True $ y=\dfrac{x^6}{6}$} 
		 {$ y=6x^5$}
	 \loigiai{
	 Ta có $\displaystyle\int{x^5}\mathrm{\,d}x=\dfrac{x^6}{6}+C$.}
\end{ex}

% nguồn: Phan1_Lop_12_Dai_so_Chuong_4
\begin{ex}%[2D4N1-2]%[ma:9b00fdd0]
	Tìm họ nguyên hàm của hàm số $f(x)=\dfrac{5+2x^4}{x^2}$. 
	 \choice
		{$\displaystyle\int  f(x)\mathrm{~d} x=\dfrac{2x^3}{3}+\dfrac{5}{x}+C$}
		{\True $\displaystyle\int  f(x)\mathrm{~d} x=\dfrac{2x^3}{3}-\dfrac{5}{x}+C$}
		 {$\displaystyle\int f(x)\mathrm{~d} x=\dfrac{2x^3}{3}+5\ln x^2+C$}
		{$\displaystyle\int f(x)\mathrm{~d}  x=2x^3-\dfrac{5}{x}+C$}
	\loigiai{ 
		Ta có $\displaystyle\int  f(x)\mathrm{~d} x=\displaystyle\int\left(2x^2+\dfrac{5}{x^2}\right)\mathrm{~d} x=\dfrac{2x^3}{3}-\dfrac{5}{x}+C$.
	} 
\end{ex}

% nguồn: Phan1_Lop_12_Dai_so_Chuong_4
\begin{ex}%[2D4N1-2]%[ma:12218c09]
	Hàm số nào sau  đây \textbf{không} là một nguyên hàm của hàm số $ y=x^3$?
	\choice 
		{$  y=\dfrac{x^4}{4}+1$}
		{$  y=\dfrac{x^4}{4}+2$}
		{$ y=\dfrac{x^4}{4}+3$} 
		{\True  $ y=3x^2$}
	 \loigiai{
	 Ta có $\displaystyle\int{x^3}\mathrm{\,d}x=\dfrac{x^4}{4}+C$.}
\end{ex}

% nguồn: Phan1_Lop_12_Dai_so_Chuong_4
\begin{ex}%[2D4N1-2]%[ma:29c7ae04]
	Nguyên hàm của hàm số $f(x)=3  x^{2}+\dfrac{1}{x}$ là
	\choice 
		{$F(x)=6 x+\ln  x+C$}
		 {$F(x)=x^{3}+\ln x+C$}
		{$F(x)=3  x^{3}+\ln |x|+C$}
		 {\True $F(x)=x^{3}+\ln |x|+C$}
	\loigiai{ 
		Ta có $F(x)=\displaystyle\int\left(3  x^{2}+\dfrac{1}{x}\right) \mathrm{\,d} x=x^{3}+\ln |x|+C$.
	} 
\end{ex}

% nguồn: Phan1_Lop_12_Dai_so_Chuong_4
\begin{ex}%[2D4N1-2]%[ma:014ed386]
	Họ nguyên hàm của hàm số $f(x)=3x^2+1$  là
	 \choice
		 {$x^3+C$}
		{$\dfrac{x^3}{3}+x+C$} 
		 {$6x+C$}
		{\True $x^3+x+C$} 
	 \loigiai{
	 Ta có $\displaystyle\int\left(3x^2+1\right) \mathrm{\,d}x=3\cdot \dfrac{x^3}{3}+x+C=x^3+x+C$.}
\end{ex}

% nguồn: Phan1_Lop_12_Dai_so_Chuong_4
\begin{ex}%[2D4N1-2]%[ma:138a67a4]
	Biết $F(x)=x^3+\dfrac{1}{x}+1$ là một nguyên hàm của hàm số $f(x)$ trên miền $(0;+\infty)$. Khẳng định nào sau  đây là đúng?
	 \choice
		 {$f(x)=\dfrac{1}{4} x^4+\ln x$}
		 {\True $f(x)=3x^2-\dfrac{1}{x^2}$}
		 {$f(x)=3x^2+\dfrac{1}{x^2}$}
		 {$f(x)=\dfrac{1}{4} x^4+\ln x+x+C$}
	 \loigiai{
		Ta có $F(x)=x^3+\dfrac{1}{x}+1$ là một nguyên hàm của hàm số $f(x)$ trên miền $(0;+\infty)$ nên $F'(x)=f(x)$ với mọi $x \in(0;+\infty)$  mà $F'(x)=\left(x^3+\dfrac{1}{x}+1\right)'=3x^2-\dfrac{1}{x^2}$.
	 }
\end{ex}

% nguồn: Phan1_Lop_12_Dai_so_Chuong_4
\begin{ex}%[2D4N1-2]%[ma:1010a398]
	Phát biểu nào sau đây  là đúng?
	\choice 
		{$\displaystyle\int \pi^3\mathrm{\,d}x  = \dfrac{\pi^4}{4}+C$}
		{$\displaystyle\int  \pi^3 \mathrm{\,d}x = 3\pi^2+C$}
		{$\displaystyle\int  \pi^3 \mathrm{\,d}x = 3\pi^2$}
		{\True $\displaystyle\int  \pi^3 \mathrm{\,d}x = \pi^3x+C$}
	\loigiai{ 
		Ta  có $\displaystyle\int \pi^3\mathrm{\,d}x=\pi^3x+C$.
	 }
\end{ex}

% nguồn: Phan1_Lop_12_Dai_so_Chuong_4
\begin{ex}%[2D4N1-2]%[ma:489e47c6]
	Cho hàm số $f(x)=6 x^{5}$,  khẳng định nào dưới đây đúng?
	\choice 
		{$\displaystyle\int  f(x) \mathrm{\,d} x=\dfrac{1}{6} x^{6}+C$}
		 {$\displaystyle\int f(x) \mathrm{\,d} x=\dfrac{1}{5} x^{6}+C$}
		{\True $\displaystyle\int f(x) \mathrm{\,d}  x=x^{6}+C$}
		{$\displaystyle\int f(x) \mathrm{\,d} x=x^{5}+C$} 
	 \loigiai{
		Ta có $\displaystyle\int 6 x^{5} \mathrm{\,d}  x=x^{6}+C$.
	} 
\end{ex}

% nguồn: Phan1_Lop_12_Dai_so_Chuong_4
\begin{ex}%[2D4N1-2]%[ma:73fe0c23]
	Với giá trị thực nào của tham số $m$ để hàm số $F(x)=m x^3+(3m+2) x^2-4x+3$ là một  nguyên hàm của hàm số $f(x)=3x^2+10x-4$?
	\choice 
		{$m=0$} 
		{\True  $m=1$}
		{$m=-1$} 
		{$m=2$} 
	 \loigiai{
		Vì $F(x)$ là một nguyên hàm của hàm số  $f(x)$ nên ta có
		$F'(x)=f(x)$.\\ 
		Khi đó  $3m x^3+2(3m+2) x-4=3x^2+10x-4\Leftrightarrow\heva{&3m=3\\
		&2(3m+2)=10}  \Leftrightarrow m=1$.
	} 
\end{ex}

% nguồn: Phan1_Lop_12_Dai_so_Chuong_4
\begin{ex}%[2D4N1-2]%[ma:f261085d]
	Tìm hàm số $y=f(x)$, biết $f'(x)=3\sqrt{x}+\dfrac{2}{\sqrt[3]x}$,  $(x > 0)$ và $f(1)=1$.
	 \choice
		 {\True $f(x)=2x \sqrt{x}+3\sqrt[3]{x^2}-4$}
		{$f(x)=2x \sqrt{x}-3\sqrt[3]{x^2}+2$} 
		 {$f(x)=\sqrt{x}+3\sqrt[3]{x^2}-3$}
		{$f(x)=x \sqrt{x}+3\sqrt[3]x-3$} 
	 \loigiai{
		Xét 
		\begin{eqnarray*}
			\displaystyle\int\left(3\sqrt{x}+\dfrac{2}{\sqrt[3]x}\right)\mathrm{~d} x&= &\displaystyle\int\left(3\sqrt{x}+\dfrac{2}{\sqrt[3]x}\right)\mathrm{~d}  x \\
			&=&  \displaystyle\int\left(3\sqrt{x}+\dfrac{2}{\sqrt[3]x}\right)\mathrm{~d} x+C.
		\end{eqnarray*}
	} 
\end{ex}

% nguồn: Phan1_Lop_12_Dai_so_Chuong_4
\begin{ex}%[2D4N1-2]%[ma:b3c1b491]
	Tìm nguyên hàm $F(x)=\displaystyle\int  \pi^2\mathrm{\,d}x$
	\choice 
		{\True $F(x)=\pi^2x+C$} 
		{$F(x)=2\pi x+C$} 
		 {$\dfrac{\pi^3}{3}+C$}
		 {$\dfrac{\pi^2x^2}{2}+C$}
	\loigiai{ 
		Ta có $F(x)=\displaystyle\int \pi^2\mathrm{\,d}x=\pi^2x+C$. 
	 }
\end{ex}

% nguồn: Phan1_Lop_12_Dai_so_Chuong_4
\begin{ex}%[2D4N1-2]%[ma:da01c509]
	Họ  nguyên hàm của hàm số $f(x)=(x+1)(x+2)$ là
	\choice 
		 {$2x+3+C$}
		{$\dfrac{x^3}{3}-\dfrac{2}{3} x^2+2x+C$} 
		{$\dfrac{x^3}{3}+\dfrac{2}{3} x^2+2x+C$} 
		{\True $\dfrac{x^3}{3}+\dfrac{3}{2}  x^2+2x+C$}
	\loigiai{ 
		Ta có $\displaystyle\int f(x)\mathrm{~d} x=\displaystyle\int(x+1)(x+2)\mathrm{~d} x=\displaystyle\int\left(x^2+3x+2\right)\mathrm{~d} x=\dfrac{x^3}{3}+\dfrac{3}{2}  x^2+2x+C$.
	 }
\end{ex}

% nguồn: Phan1_Lop_12_Dai_so_Chuong_4
\begin{ex}%[2D4N1-2]%[ma:bf03a4a9]
	 Họ nguyên hàm của hàm số $f(x)=x^6$ là
	 \choice
		 {$x^7+C$}
		{\True $\dfrac{x^7}{7}+C$} 
		 {$6x^5+C$}
		 {$x^5+C$}
	 \loigiai{
		Ta có $\displaystyle\int f(x)\mathrm{\,d}x =  \int x^6\mathrm{\,d}x = \dfrac{x^7}{7}+C$.\\
		Do đó, họ nguyên hàm của hàm số $f(x)=x^6$  là $\dfrac{x^7}{7}+C$.
	} 
\end{ex}

% nguồn: Phan1_Lop_12_Dai_so_Chuong_4
\begin{ex}%[2D4N1-2]%[ma:7a8ab21f]
	 Tìm một nguyên hàm $F(x)$ của hàm số $f(x)=2x+3\sqrt{x}$ thoả mãn $F(1)=0$.
	\choice 
		 {$F(x)=x^2+3\sqrt{x^3}$}
		 {$F(x)=x^2+2\sqrt[3]{x^2}$}
		{$F(x)=x^2+3\sqrt[3]{x^2}-4$} 
		{\True  $F(x)=x^2+2\sqrt{x^3}-3$}
	 \loigiai{
		 Ta có $\displaystyle\int(2x+3\sqrt{x})\mathrm{~d} x=x^2+2x \sqrt{x}+C$ mà $F(1)=0\Rightarrow C=-3$.\\
		Vậy  $F(x)=x^2+2x \sqrt{x}-3$.
	} 
\end{ex}

% nguồn: Phan1_Lop_12_Dai_so_Chuong_4
\begin{ex}%[2D4N1-2]%[ma:7f363310]
	Nguyên hàm của hàm  số $f(x)=x^2+1$ là
	\choice 
		{$x^3+x+C$} 
		 {$3x^3+x+C$}
		 {$x^3+C$}
		{\True $\dfrac{1}{3}x^3+x+C$} 
	\loigiai{ 
		Ta có $\displaystyle\int  (x^2+1) \mathrm{\,d}x=\dfrac{1}{3}x^3+x+C$.
	} 
\end{ex}

% nguồn: Phan1_Lop_12_Dai_so_Chuong_4
\begin{ex}%[2D4N1-2]%[ma:37db3c8d]
	 Họ nguyên hàm của hàm số $f(x)=2x+1$ là
	 \choice
		{$x^2+1+C$} 
		{\True $x^2+x+C$} 
		{$x^2  + C$}
		{$2x+C$} 
	\loigiai{ 
		$\displaystyle\int (2x+1)dx = x^2+x+C$ 
	} 
\end{ex}

% nguồn: Phan1_Lop_12_Dai_so_Chuong_4
\begin{ex}%[2D4N1-2]%[ma:7d81f9c0]
	 Họ các nguyên hàm của hàm số $f(x)=3x^2+1$ là
	 \choice
		 {$x^3+C$}
		 {$\dfrac{x^3}{3}+x+C$}
		 {\True $x^3+x+C$}
		 {$6x+C$}
	 \loigiai{
		Ta có $\displaystyle\int f(x)\mathrm{~d} x=\displaystyle\int\left(3x^2+1\right)\mathrm{~d} x=3\cdot \dfrac{x^3}{3}+x+C=x^3+x+C$. 
	 }
\end{ex}

% nguồn: Phan1_Lop_12_Dai_so_Chuong_4
\begin{ex}%[2D4N1-2]%[ma:0f4d1ac9]
	Họ nguyên hàm  của hàm số $f(x)=3\sqrt{x}+x^{2024}$ là
	\choice 
		{$\sqrt{x}+\dfrac{x^{2023}}{673}+C$} 
		{\True $2\sqrt{x^3}+\dfrac{x^{2025}}{2025}+C$} 
		 {$\dfrac{1}{\sqrt{x}}+\dfrac{x^{2024}}{2025}+C$}
		 {$\dfrac{1}{2\sqrt{x}}+6054x^{2025}+C$}
	 \loigiai{
		Ta  có
		\begin{eqnarray*}
			 \displaystyle\int\left(3\sqrt{x}+x^{2024}\right)\mathrm{~d} x& =& \\
			&=&\displaystyle\int\left(3x^{\tfrac{1}{2}}+x^{2024}\right)\mathrm{~d} x\\ 
			&=&3\cdot  \dfrac{x^{\tfrac{3}{2}}}{\dfrac{3}{2}}+\dfrac{x^{2025}}{2025}+C\\
			 &=&2\sqrt{x^3}+\dfrac{x^{2025}}{2025}+C.
		\end{eqnarray*}
	 }
\end{ex}

% nguồn: Phan1_Lop_12_Dai_so_Chuong_4
\begin{ex}%[2D4N1-2]%[ma:92a0d05f]
	$\displaystyle\int x^4  \mathrm{\,d} x$ bằng
	 \choice
		 {\True $\dfrac{1}{5} x^5+C$}
		{$4x^3+C$} 
		{$x^5+C$} 
		{$5x^5+C$} 
	\loigiai{ 
		$\displaystyle\int x^4  \mathrm{\,d} x=\dfrac{1}{5} x^5+C$.
	} 
\end{ex}

% nguồn: Phan1_Lop_12_Dai_so_Chuong_4
\begin{ex}%[2D4N1-2]%[ma:a0a4dc1b]
	Nguyên hàm  của hàm số $f(x)=4x^3+2x$ là
	 \choice
		 {$\dfrac{1}{4}x^4+\dfrac{1}{2}x^2+C$}
		 {$3x^2+1+C$}
		{\True  $x^4+x^2+C$}
		 {$x^3+x+C$}
	 \loigiai{
		Áp  dụng công thức.
	} 
\end{ex}

% nguồn: Phan1_Lop_12_Dai_so_Chuong_4
\begin{ex}%[2D4H1-2]%[ma:fc5d75b5]
	Tìm một nguyên hàm $F(x)$ của  hàm số $f(x)=2x-\mathrm{e}^x$, biết $F(0)=-2$.
	\choice 
		 {$F(x)=x-\mathrm{e}^x+1$}
		{\True $F(x)=x^2-\mathrm{e}^x-1$} 
		 {$F(x)=x^2-\mathrm{e}^x-2$}
		{$F(x)=x^2-\mathrm{e}^x+1$} 
	 \loigiai{
		 $F(x)=\displaystyle\int (2x-\mathrm{e}^x)\mathrm{\,d}x=x^2-\mathrm{e}^x+C$.\\
		Ta có $F(0)=-2\Rightarrow-1+C=-2\Leftrightarrow C=-1\Rightarrow  F(x)=x^2-\mathrm{e}^x-1$.
	 }
\end{ex}

% nguồn: Phan1_Lop_12_Dai_so_Chuong_4
\begin{ex}%[2D4H1-2]%[ma:81375bb8]
	Gọi $F(x)$ là một nguyên hàm của hàm số $f(x)=(2x-3)^2$ thỏa $F(0)=\dfrac{1}{3}$. Tính giá trị của biểu  thức $T=\log_2[3F(1)-2F(2)]$.
	 \choice
		{$T=10$} 
		{$T=4$} 
		 {\True $T=2$}
		{$T=-4$} 
	 \loigiai{
		Ta có  $F(x)=\displaystyle\int(2x-3)^2\mathrm{~d} x=\dfrac{1}{2} \displaystyle\int(2x-3)^2\mathrm{~d}(2x-3)=\dfrac{1}{2} \cdot \dfrac{(2x-3)^3}{3}+C$.\\
		Ta có $F(0)=\dfrac{1}{3} \Leftrightarrow \dfrac{1}{2} \cdot  \dfrac{(0-3)^3}{3}+C=\dfrac{1}{3} \Leftrightarrow C=\dfrac{29}{6}$.\\
		Do đó $F(x)=\dfrac{1}{2} \cdot \dfrac{(2x-3)^3}{3}+\dfrac{29}{6}$ nên $F(1)=\dfrac{1}{2} \cdot-\dfrac{1}{3}+\dfrac{29}{6}=\dfrac{14}{3}; F(2)=\dfrac{1}{2} \cdot \dfrac{1}{3}+\dfrac{29}{6}=5$.\\ 
		 $\Rightarrow T=\log_2[3F(1)-2F(2)]=\log_2\left(3\cdot \dfrac{14}{3}-2\cdot5\right)=\log_2 4=2$.
	 }
\end{ex}

% nguồn: Phan1_Lop_12_Dai_so_Chuong_4
\begin{ex}%[2D4H1-2]%[ma:120aa15a]
	Cho hàm số $y=f(x)$ xác định trên $\mathbb{R}$ thỏa mãn $f'(x)=x^2-2x-3$ và $f(0)=1$. Khi đó $f(x)$  bằng
	\choice 
		{$\dfrac{x^3}{3}-x^2-3x$} 
		{$\dfrac{x^3}{3}-x^2-3x+3$} 
		{\True  $\dfrac{x^3}{3}-x^2-3x+1$}
		 {$\dfrac{x^3}{3}-x^2-3x+2$}
	\loigiai{Ta có $f'(x)=x^2-2x-3$. \\ 
		Suy ra  $f(x)=\displaystyle\int f'(x)\,\mathrm{d}x=\displaystyle\int(x^2-2x-3)\,\mathrm{d}x=\dfrac{x^3}{3}-x^2-3x+C$.\\
		Lại có  $f(0)=1\Rightarrow C=1$.\\
	Vậy $f(x)=\dfrac{x^3}{3}-x^2-3x+1$.} 
\end{ex}

% nguồn: Phan1_Lop_12_Dai_so_Chuong_4
\begin{ex}%[2D4H1-2]%[ma:222c157c]
	Giá  trị của $\displaystyle\int_{0}^{2}|x^2-x|\mathrm{\,d}x$ bằng
	\choice 
		{$\dfrac{2}{3}$} 
		{\True $1$} 
		{$\dfrac{1}{3}$} 
		{$2$} 
	\loigiai{ 
		Ta  có: $x^2-x=0 \Leftrightarrow x=0$ hoặc $x=1$.\\
		Do đó,\begin{align*} \displaystyle\int_{0}^{2}|x^2-x|\mathrm{\,d}x&=\displaystyle\int_{0}^{1}|x^2-x|\mathrm{\,d}x+\displaystyle\int_{1}^{2}|x^2-x|\mathrm{\,d}x\\&=\left|\displaystyle\int_{0}^{1}(x^2-x)\mathrm{\,d}x\right|+\left|\displaystyle\int_{1}^{2}(x^2-x)\mathrm{\,d}x\right|\\&=\left|\left(\dfrac{x^3}{3}-\dfrac{x^2}{2}\right)\bigg|_{0}^{1}\right|+\left|\left(\dfrac{x^3}{3}-\dfrac{x^2}{2}\right)\bigg|_{1}^{2}\right|=\left|-\dfrac{1}{6}\right|+\left|\dfrac{5}{6}\right|=1
	\end{align*}
 }
\end{ex}

% nguồn: Phan1_Lop_12_Dai_so_Chuong_4
\begin{ex}%[2D4H1-2]%[ma:c3561220]
	Biết $F(x)=x^3+\dfrac{1}{x}+1$ là một nguyên hàm của hàm số $f(x)$ trên khoảng $(0;+\infty)$. Khẳng định  nào sau đây là đúng?
	\choice 
		 {$f(x)=\dfrac{1}{4}x^4+\ln x$}
		 {\True $f(x)=3x^2-\dfrac{1}{x^2}$}
		{$f(x)=3x^2+\dfrac{1}{x^2}$} 
		{$f(x)=\dfrac{1}{4}x^4+\ln  x+x+C$}
	\loigiai{Ta có  $F(x)=x^3+\dfrac{1}{x}+1$.
	 Suy ra $f(x)=F'(x)=3x^2-\dfrac{1}{x^2}$.}
\end{ex}

% nguồn: Phan1_Lop_12_Dai_so_Chuong_4
\begin{ex}%[2D4H1-2]%[ma:afd461d9]
	Cho hàm  số $f(x)=\heva{&2x+3\,\text{ khi } x \geq 1\\
	&3x^2+2\,\text{ khi } x < 1}$. Giả sử $F$ là nguyên hàm của $f$ trên $\mathbb{R}$  thỏa mãn $F(0)=2$. Giá trị $F(-1)+2F(2)$ bằng
	 \choice
		 {$23$}
		{$11$} 
		 {$10$}
		{\True  $21$}
	 \loigiai{
		Do $F$ là nguyên hàm  của $f$ trên $\mathbb{R}$ nên
		$$F(x)=\heva{&F_1(x)=\displaystyle\int(2x+3)\mathrm{~d} x=x^2+3x+C_1& \text{ khi  } x \geq 1\\ &F_2(x)=\displaystyle\int\left(3x^2+2\right)\mathrm{~d} x=x^3+2x+C_2& \text{ khi } x < 1.}$$
		Theo bài ra, ta có $F(x)$ có đạo  hàm trên R nên $F(x)$ liên tục trên đó và $F(0)=2$.\\
		Điều này tương  đương với $\heva{&F_1(1)=F_2(1) \\
		 &F_2(0)=2} \Leftrightarrow\heva{&4+C_1=3+C_2\\
		 &C_2=2} \Leftrightarrow\heva{&C_1=1\\
		&C_2=2.}$\\ 
		Do đó, $F_1(x)=x^2+3x+1, F_2(x)=x^3+2x+2$,  suy ra $F(-1)=F_2(-1)=-1, F(2)=F_1(2)=11$.\\ Vậy $F(-1)+2F(2)=-1+2\cdot 11=21$
	} 
\end{ex}

% nguồn: Phan1_Lop_12_Dai_so_Chuong_4
\begin{ex}%[2D4H1-2]%[ma:93e9fff5]
	Cho $a$ là số dương khác 1. Hàm số $y = \log_a x$ là một nguyên hàm  của hàm số
	\choice 
		{$y = \dfrac{1}{x \ln\mathrm{e}}$} 
		{\True $y = \dfrac{1}{x \ln a}$} 
		{$y  = \dfrac{1}{x}$}
		{$y = \dfrac{\ln a}{x}$} 
	 \loigiai{
		Ta có $y'  = (\log_a x)'=\dfrac{1}{x \ln a}$.
	} 
\end{ex}

% nguồn: Phan1_Lop_12_Dai_so_Chuong_4
\begin{ex}%[2D4H1-2]%[ma:dbe7caf0]
	Biết hàm số $F(x)$ là một nguyên hàm  của hàm số $f(x)=3x^2-6x+4$ trên và $F(1)=7$. Tính $F(-1)$.
	 \choice
		{$1$} 
		{$2$} 
		{\True $-3$} 
		{$-1$} 
	 \loigiai{
		 Ta có $F(x)=\displaystyle\int \left(3x^2-6x+4\right) \mathrm{\,d}x=x^3-3x^2+4x+C$.\\
		Do $F(1)=7 \Leftrightarrow 1^3-3\cdot  1^2+4\cdot 1 +C=7 \Leftrightarrow C=5$.\\
		 Vậy $F(x)=x^3-3x^2+4x+5$.\\
		Khi  đó $F(-1)=(-1)^3-3\cdot (-1)^2+4\cdot (-1)+5=-3$.
	} 
\end{ex}

% nguồn: Phan1_Lop_12_Dai_so_Chuong_4
\begin{ex}%[2D4H1-2]%[ma:9f206a3a]
	 Biết rằng $\displaystyle\int_{0}^{2}f(x)\mathrm{\,d}x=-4$. Giá trị của $\displaystyle\int_{0}^{2}[3x-2f(x)]\mathrm{\,d}x$ bằng
	\choice 
		{$-2$} 
		 {$12$}
		{\True $14$} 
		{$22$} 
	\loigiai{ 
		 Ta có $\displaystyle\int_{0}^{2}[3x-2f(x)]\mathrm{\,d}x=\displaystyle\int_{0}^{2}3x\mathrm{\,d}x-2\displaystyle\int_{0}^{2}f(x)\mathrm{\,d}x=\left(\dfrac{3x^2}{2}\right)\bigg|_{0}^{2}-2(-4)=6+8=14$
	 }
\end{ex}

% nguồn: Phan1_Lop_12_Dai_so_Chuong_4
\begin{ex}%[2D4H1-2]%[ma:68fd0c15]
	Cho hàm số $f(x)=\heva{&2x+5\,\text{ khi } x  \geq 1\\
	&3x^2+4\, \text{ khi } x < 1}$. Giả sử $F$ là nguyên hàm của $f$ trên $\mathbb{R}$  thỏa mãn $F(0)=2$. Giá trị của $F(-1)+2F(2)$ bằng
	 \choice
		 {\True $27$}
		{$29$} 
		 {$12$}
		{$33$} 
	 \loigiai{
		Theo giả thiết $F$  là một nguyên hàm của $f(x)$ trên $\mathbb{R}$ nên ta có
		$$f(x)=\heva{&2x+5\text{ khi } x \geq  1\\
		&3x^2+4\text{ khi } x <  1} \Rightarrow F(x)=\heva{&x^2+5x+C_1& \text{ khi } x \geq 1\\ &x^3+4x+C_2& \text{ khi } x < 1.}$$
		Vì $F(0)=0\Rightarrow C_2=2$.\\ 
		Mặt khác, $F(x)$ liên tục trên R nên liên tục tại $x=1$ nên ta  có
		$$\lim\limits_{x \to  1^{+}} F(x)=\lim\limits_{x \to 1^{-}} F(x) \Leftrightarrow 6+C_1=5+C_2 \Rightarrow C_1=C_2-1=1.$$
		Vậy $F(x)=\heva{&x^2+5x+1\text{ khi }  x \geq 1\\
		&x^3+4x+2\text{ khi } x < 1} \Rightarrow  F(-1)+2F(2)=-3+2\cdot15=27$.
	} 
\end{ex}

% nguồn: Phan1_Lop_12_Dai_so_Chuong_4
\begin{ex}%[2D4H1-2]%[ma:a5a1850c]
	Nguyên hàm $F(x)$ của  hàm số $f(x)=3x^2+4x$ thỏa mãn $F(1)=5$ là
	\choice 
		 {$F(x)=9x^3+8x^2-12$}
		{$F(x)=x^3+2x^2+5$} 
		 {$F(x)=6x+4$}
		{\True  $F(x)=x^3+2x^2+2$}
	 \loigiai{Ta có $F(x)=\displaystyle\int f(x)\mathrm{\,d}x=\displaystyle\int\left(3x^2+4x\right)\mathrm{\,d}x=x^3+2x^2+C$.\\
		Do  $F(1)=5$ nên $C+3=5\Leftrightarrow C=2$.\\
	Vậy $F(x)=x^3+2x^2+2$.} 
\end{ex}

% nguồn: Phan1_Lop_12_Dai_so_Chuong_4
\begin{ex}%[2D4H1-2]%[ma:76278cb3]
	 Biết $F(x)$ là một nguyên hàm của hàm số $f(x)$ trên đoạn $[a;b]$ và $2F(a)-7=2F(b)$. Tính tích phân $I=\displaystyle\int_{a}^{b} f(x)\mathrm{\,d}x$.
	 \choice
		 {$I=-2$}{$I=\dfrac{7}{2}$}{\True $I=-\dfrac{7}{2}$}{$I=2$}
	\loigiai{Ta có: $2F(a)-7=2F(b)\Leftrightarrow  F(b)-F(a)=-\dfrac{7}{2}$.\\
		$I=\displaystyle\int_{a}^{b} f(x)\mathrm{\,d}x=F(b)-F(a)=-\dfrac{7}{2}$. 
	} 
\end{ex}


%%%%%%%%%%%%%%%% ĐÚNG/SAI %%%%%%%%%%%%%%%%

% nguồn: Phan1_Lop_12_Dai_so_Chuong_4
\begin{ex}%[2D4N1-2]%[ma:7b38bf99]
	Cho hàm số $f(x)=4x^3-6x$.  Biết $F(x)$ là một nguyên hàm của $f(x)$ và $F(0)=2$.
	 \choiceTF
		{$F(x)=f'(x)$} 
		 {\True $F'(x)=f(x)$}
		{\True $F(x)=x^4-3x^2+2$} 
		{$F(1)=3$} 
	 \loigiai{
		\begin{itemchoice}
			\itemch Theo định nghĩa $F(x)$  là một nguyên hàm của $f(x)$ nên $F'(x)=f(x)$.
			\itemch $F'(x)=f(x)$. 
			\itemch Ta có $F(x)=\displaystyle\int 4x^3\mathrm{~d}  x-\displaystyle\int 6x \mathrm{~d} x=x^4-3x^2+C$.\\
			Theo giả thiết $F(0)=2$ suy ra  $C=2\Rightarrow F(x)=x^4-3x^2+2$.
			\itemch $F(1)=1^4-3\cdot1^2+2=0$. 
		\end{itemchoice}
	} 
\end{ex}

% nguồn: Phan1_Lop_12_Dai_so_Chuong_4
\begin{ex}%[2D4H1-2]%[ma:f8b55fe4]
	Cho hàm số $f(x)=2x+m$  và $G(x)=x^3+m x^2+3x+m$ với $m \in \mathbb{R}$. Gọi $F(x)$ làm một nguyên hàm của $f(x)$ sao cho $F(0)=1$.
	\choiceTF 
		{Khi $m=3$  thì $\displaystyle\int f(x)\mathrm{~d} x=x^2-3x+C$}
		{Khi $m=2$ thì $G(x)$ là một nguyên  hàm của $f(x)$}
		{\True Khi $m=-1$ thì $\displaystyle\int(G(x)-f(x))\mathrm{~d} x=\dfrac{1}{4} x^4-\dfrac{1}{3} x^3+\dfrac{1}{2}  x^2+C$}
		{\True Có $2$ giá trị nguyên dương của  m để $F(x) \geq 0, \forall x \in \mathbb{R}$}
	 \loigiai{
		\begin{itemchoice}
			\itemch Khi $m=3$  thì $\displaystyle\int f(x)\mathrm{~d} x=\displaystyle\int(2x+3)\mathrm{~d} x=\displaystyle\int 2x \mathrm{~d} x+\displaystyle\int 3\mathrm{~d} x=x^2+3x+C$.
			\itemch Khi $m=2$,  $f(x)=2x+2$; $G(x)=x^3+2x^2+3x+2$.\\
			Ta có $G'(x)=3x^2+4x+3\Rightarrow G'(1)=10\neq f(1)$ nên khi $m=2$, $G(x)$ là không phải là một nguyên hàm  của $f(x)$.
			\itemch Khi  $m=-1$, $f(x)=2x-1$; $G(x)=x^3-x^2+3x-1$.
			\begin{eqnarray*}
				\displaystyle\int(G(x)-f(x))\mathrm{~d} x &=&\displaystyle\int\left(x^3-x^2+x\right)\mathrm{~d} x  \\
				&=& \displaystyle\int x^3\mathrm{~d} x-\displaystyle\int x^2\mathrm{~d} x+\displaystyle\int  x \mathrm{~d} x\\
				&=&\dfrac{1}{4} x^4-\dfrac{1}{3} x^3+\dfrac{1}{2} x^2+C. 
			\end{eqnarray*}
			\itemch Ta có $F(x)=\displaystyle\int f(x)\mathrm{~d} x=\displaystyle\int 2x \mathrm{~d} x+\displaystyle\int m  \mathrm{~d} x=x^2+m x+C$.\\
			 Ta có $F(0)=1$ suy ra $C=1$.\\
			Khi  đó $F(x)=x^2+m x+1$.\\
			$F(x) \geq 0, \forall x \in \mathbb{R} \Leftrightarrow \Delta=m^2-4\leq 0\Leftrightarrow-2\leq m  \leq 2$.\\
			Vì $m$ nguyên  dương nên $m \in\{1; 2\}$.
		\end{itemchoice}
	} 
\end{ex}

% nguồn: Phan1_Lop_12_Dai_so_Chuong_4
\begin{ex}%[2D4H1-2]%[ma:9854f712]
	Cho hàm số $F(x)=\displaystyle\int  \left(2x+\sqrt{x}\right)\mathrm{\,d}x$ (với $x>0$). Khi đó
	\choiceTF 
		 {\True $F(x)=2\displaystyle\int x\mathrm{\,d}x+\displaystyle\int \sqrt{x}\mathrm{\,d}x+C$ với $C\in \mathbb{R}$}
		 {\True $G(x)=F(x)+2024 \Rightarrow G(x)=\displaystyle\int \left(2x+\sqrt{x}\right)\mathrm{\,d}x$}
		{$F(x)=x+\dfrac{2}{3}x\sqrt{x}+C$} 
		{$F(1)=\dfrac{2}{3} \Rightarrow  F(4)=\dfrac{28}{3}$}
	 \loigiai{
		\begin{itemchoice}
			\itemch Vì  $\displaystyle\int kf(x)\mathrm{\,d}x = k\displaystyle\int f(x)\mathrm{\,d}x$, $k\ne 0$; $\displaystyle\int \left(f(x)+g(x)\right)\mathrm{\,d}x = \displaystyle\int f(x)\mathrm{\,d}x + \displaystyle\int g(x)\mathrm{\,d}x$.
			\itemch Ta có $F(x)$, $G(x)$ là  nguyên hàm của hàm số $f(x)$ thì $F(x)=G(x)+C$.
			\itemch Ta có $F(x)=\displaystyle\int \left(2x+\sqrt{x}\right)\mathrm{\,d}x = \displaystyle\int \left(2x+x^{\frac{1}{2}}\right)\mathrm{\,d}x = x^2+\dfrac{2}{3}x\sqrt{x}+C$ suy ra khẳng định sai. 
			\itemch Ta có $F(1)=1+\dfrac{2}{3}+C=\dfrac{2}{3} \Rightarrow C=-1 \Rightarrow F(x)=x^2+\dfrac{2}{3}x\sqrt{x}-1 \Rightarrow  F(4)=\dfrac{25}{3}$.
		\end{itemchoice}
	} 
\end{ex}

% nguồn: Phan1_Lop_12_Dai_so_Chuong_4
\begin{ex}%[2D4H1-2]%[ma:3b705257]
	Cho $F(x)$ là một nguyên hàm của hàm số $f(x)=2x-2$, biết  rằng $F(1)=1$. Xét tính đúng sai của các khẳng định sau
	 \choiceTF
		{\True $F(x)$ luôn xác định trên  $\mathbb{R}$}
		{$F(x)=x^2-2x$} 
		{\True $F(x) > 0$, với  mọi $x \in \mathbb{R}$}
		 {$F(-2)=2$}
	\loigiai{\begin{itemchoice}
		\itemch \textbf{Đúng}.\\ 
		$F(x)=x^2-2x+C$. 
		Suy ra $F(x)$ luôn xác định trên $\mathbb{R}$. 
		 \itemch \textbf{Sai}.\\
		$F(x)=x^2-2x+C$ mà $F(1)=1 \Rightarrow  C=2$.\\
		Vậy $F(x)=x^2-2x+2$. 
		 \itemch \textbf{Đúng}.\\
		Ta có  $F(x)=x^2-2x+2=(x-1)^2+1>0$, $\forall x\in \mathbb{R}$.
		\itemch  \textbf{Sai}.\\
		 $F(-2)=(-2)^2-2\cdot(-2)+2=10$.
\end{itemchoice}}
\end{ex}

% nguồn: Phan1_Lop_12_Dai_so_Chuong_4
\begin{ex}%[2D4H1-2]%[ma:b4a94fff]
	Cho hàm số $f(x)=\dfrac{2x^2+3x+2}{x+1}$ với  $x\ne -1$. Khi đó
	 \choiceTF
		{\True $f(x)=2x+1+\dfrac{1}{x+1}$} 
		{$\displaystyle\int f(x)\mathrm{\,d}x =  x^2+x+\ln{(x+1)}+C$}
		{\True Nguyên hàm $F(x)$ của hàm số  $f(x)=\dfrac{2x^2+3x+2}{x+1}$ thỏa mãn $F(-2)=3$ là $F(x)=x^2+x+\ln{\left|x+1\right|}+1$}
		 {Bất phương trình $F(x)<x^2+x+2$ có tập nghiệm là $T=\left(-e-1;e-1\right)$}
	\loigiai{ 
		\begin{itemchoice}
			\itemch Ta  có $f(x)=2x+1+\dfrac{1}{x+1}$.
			\itemch Ta có $\displaystyle\int f(x)\mathrm{\,d}x = \displaystyle\int \left(2x+1+\dfrac{1}{x+1}\right)\mathrm{\,d}x = x^2+x+\ln{\left|x+1\right|}+C$. 
			\itemch  Ta có $\displaystyle\int f(x)\mathrm{\,d}x = \displaystyle\int \left(2x+1+\dfrac{1}{x+1}\right)\mathrm{\,d}x = x^2+x+\ln{\left|x+1\right|}+C$. \\
			Ta có $$F(-2)=3 \Leftrightarrow (-2)^2+(-2)+\ln{\left|-2+1\right|}+C=3 \Leftrightarrow  C=1.$$ Vậy $F(x)=x^2+x+\ln{\left|x+1\right|}+1$.
			\itemch 
			Ta có \begin{eqnarray*}
			&&F(x)<x^2+x+2\\ 
			 &\Leftrightarrow&x^2+x+\ln{\left|x+1\right|}+1<x^2+x+2\\
			&\Leftrightarrow&\ln{\left|x+1\right|}<1\\ 
			 &\Leftrightarrow&\left|x+1\right|<e\\
			 &\Leftrightarrow&-e<x+1<e\\
			 &\Leftrightarrow&-e-1<x<e-1, x\neq -1.
		\end{eqnarray*}
		Vậy tập nghiệm  của bất phương trình $F(x)<x^2+x+2$ là $T=\left(-e-1;e-1\right)\setminus\{-1\}$.
	\end{itemchoice}
} 
\end{ex}

% nguồn: Phan1_Lop_12_Dai_so_Chuong_4
\begin{ex}%[2D4H1-2]%[ma:65e09d53]
	Cho hàm số $f(x)=x^2$. Khi  đó
	 \choiceTF
		{\True $\displaystyle\int f(x)\mathrm{\,d}x  = \dfrac{x^3}{3}+C$}
		{Gọi $F(x)$ là một  nguyên hàm của $f(x)$. Biết $F(3)=1$ thì $F(4)=\dfrac{4}{3}$}
		{\True $\displaystyle\int f(2x+1)\mathrm{\,d}x = \displaystyle\int (2x+1)^2\mathrm{\,d}x  = \dfrac{4}{3}x^3 + 2x^2 + x + C$}
		 {$\displaystyle\int \left[x f(x-2)\right]\mathrm{\,d}x = \dfrac{x^4}{4} - \dfrac{x^3}{3} + 2x^2 + C$}
	\loigiai{ 
		\begin{itemchoice}
			\itemch Ta có $\displaystyle\int f(x)\mathrm{\,d}x = \displaystyle\int  x^2\mathrm{\,d}x = \dfrac{x^3}{3}+C$.
			\itemch Ta  có $F(x)=\dfrac{x^3}{3}+C$, $F(3)=1 \Rightarrow \dfrac{27}{3}+C=1 \Rightarrow C=-8$ nên $F(x)=\dfrac{x^3}{3}-8$. \\
			Do đó  $F(4)=\dfrac{64}{3}-8=\dfrac{40}{3}$.
			\itemch Ta có $\displaystyle\int f(2x+1)\mathrm{\,d}x  = \displaystyle\int (2x+1)^2\mathrm{\,d}x = \displaystyle\int (4x^2+4x+1)\mathrm{\,d}x = \dfrac{4}{3}x^3 + 2x^2 + x + C$.
			\itemch  Ta có
			\begin{eqnarray*}
				&\displaystyle\int \left[x f(x-2)\right]\mathrm{\,d}x&= \displaystyle\int x(x-2)^2\mathrm{\,d}x\\ 
				&&= \displaystyle\int x(x^2 - 4x + 4)\mathrm{\,d}x\\ 
				&&= \displaystyle\int (x^3 - 4x^2  + 4x)\mathrm{\,d}x\\
				&&= \dfrac{x^4}{4} - \dfrac{4x^3}{3} + 2x^2 +  C.
			\end{eqnarray*}
		\end{itemchoice}
	} 
\end{ex}

% nguồn: Phan1_Lop_12_Dai_so_Chuong_4
\begin{ex}%[2D4H1-2]%[ma:618c3d89]
	Cho hàm số $f(x)=x^3-4x+5$. Gọi $F(x)$ là một nguyên  hàm của $f(x)$. Biết $F(1)=3$. Khi đó
	\choiceTF 
		{\True $\displaystyle\int  (x^3-4x+5)\mathrm{\,d}x = \dfrac{x^4}{4}-2x^2+5x+C$}
		 {$F(0)=2$}
		{$\displaystyle\int \left[f(x)+f'(x)\right]\mathrm{\,d}x =  \dfrac{x^4}{4}+x^3-2x^2+9x+C$}
		{\True $\displaystyle\int f(x+1)\mathrm{\,d}x  = \dfrac{x^4}{4}+x^3-\dfrac{1}{2}x^2+2x+C$}
	 \loigiai{
		\begin{itemchoice}
			\itemch Ta có $\displaystyle\int  (x^3-4x+5)\mathrm{\,d}x = \dfrac{x^4}{4}-2x^2+5x+C$.
			 \itemch Ta có $F(1)=3 \Leftrightarrow \dfrac{1^4}{4}-2\cdot1^2+5\cdot1+C = 3 \Leftrightarrow \dfrac{13}{4}+C = 3 \Rightarrow C=-\dfrac{1}{4}$ nên $F(0) = -\dfrac{1}{4}$.
			 \itemch Ta có
			\begin{eqnarray*}
				 &\displaystyle\int \left[f(x)+f'(x)\right]\mathrm{\,d}x&= \displaystyle\int \left(x^3-4x+5 + 3x^2-4\right)\mathrm{\,d}x\\
				&&= \displaystyle\int  (x^3+3x^2-4x+1)\mathrm{\,d}x\\
				&&= \dfrac{x^4}{4}+x^3-2x^2+x+C. 
			\end{eqnarray*}
			 \itemch
			Ta  có
			\begin{eqnarray*}
				&\displaystyle\int f(x+1)\mathrm{\,d}x&= \displaystyle\int \left[(x+1)^3 - 4(x+1) +  5\right]\mathrm{\,d}x \\
				&&= \displaystyle\int \left(x^3+3x^2+3x+1 -4x -4 +5\right)\mathrm{\,d}x\\ 
				&&= \displaystyle\int  (x^3+3x^2 -x +2)\mathrm{\,d}x\\
				&&= \dfrac{x^4}{4}+x^3 - \dfrac{1}{2}x^2  + 2x + C.
			\end{eqnarray*}
		\end{itemchoice}
	 }
\end{ex}

% nguồn: Phan1_Lop_12_Dai_so_Chuong_4
\begin{ex}%[2D4H1-2]%[ma:99285b25]
	Cho hàm số $f(x)=\dfrac{2}{x^2-1}$ với $x\ne \pm  1$. Khi đó
	 \choiceTF
		{\True $f(x)=\dfrac{1}{x-1}-\dfrac{1}{x+1}$} 
		 {$\displaystyle\int f(x)\mathrm{\,d}x = \ln{\left|\dfrac{x+1}{x-1}\right|}+C$}
		{\True Nguyên hàm  $F(x)$ của hàm số $f(x)=\dfrac{2}{x^2-1}$ thỏa mãn $F(-2)=\ln{(3e)}$ là $F(x)=\ln{\left|\dfrac{x-1}{x+1}\right|}+1$}
		{\True Phương trình $F(x)=\ln{(2\mathrm{e})}$  có 2 nghiệm $x_1$, $x_2$. Khi đó $S=x_1x_2=1$}
	\loigiai{ 
		\begin{itemchoice}
			\itemch Ta có $f(x)=\dfrac{(x+1)-(x-1)}{(x+1)(x-1)}=\dfrac{1}{x-1}-\dfrac{1}{x+1}$. 
			\itemch Ta có $\displaystyle\int f(x)\mathrm{\,d}x = \displaystyle\int \left(\dfrac{1}{x-1}-\dfrac{1}{x+1}\right)\mathrm{\,d}x =  \ln{\left|x-1\right|}-\ln{\left|x+1\right|}+C = \ln{\left|\dfrac{x-1}{x+1}\right|}+C$.
			\itemch Ta có $\displaystyle\int f(x)\mathrm{\,d}x = \displaystyle\int \left(\dfrac{1}{x-1}-\dfrac{1}{x+1}\right)\mathrm{\,d}x = \ln{\left|x-1\right|}-\ln{\left|x+1\right|}+C =  \ln{\left|\dfrac{x-1}{x+1}\right|}+C$. \\
			Ta có $F(-2)=\ln{(3\mathrm{e})} \Leftrightarrow \ln{\left|\dfrac{-2-1}{-2+1}\right|}+C = \ln{(3\mathrm{e})} \Leftrightarrow \ln{3}+1 = \ln{(3\mathrm{e})}  \Rightarrow C=1$. \\
			Vậy $F(x)=\ln{\left|\dfrac{x-1}{x+1}\right|}+1$. 
			 \itemch Ta có
			\begin{eqnarray*}
				 &&F(x)=\ln{(2\mathrm{e})}\\
				&\Leftrightarrow&\ln{\left|\dfrac{x-1}{x+1}\right|}+1 = 1+\ln{2}\\ 
				&\Leftrightarrow&\left|\dfrac{x-1}{x+1}\right|=2\\ 
				&\Leftrightarrow&\left|x-1\right|=2\left|x+1\right|\\ 
				&\Leftrightarrow&\hoac{&x=-3 \\&  x=-\dfrac{1}{3}.}
			\end{eqnarray*}
			Cả hai nghiệm trên đều thỏa mãn. Vậy $S = x_1x_2 = (-3)\cdot\left(-\dfrac{1}{3}\right) =  1$.
		\end{itemchoice}
	 }
\end{ex}

% nguồn: Phan1_Lop_12_Dai_so_Chuong_4
\begin{ex}%[2D4H1-2]%[ma:59d726a2]
	Gọi hàm số $F(x) = \dfrac{2x^2-1}{x}$ là một nguyên  hàm của hàm số $f(x)$, với $x \neq 0$.
	 \choiceTF
		{$f(x) = 2 - \dfrac{1}{x^2}$} 
		{\True  $\displaystyle\int f(x) \mathrm{\,d}x = \dfrac{2x^2-1}{x} + C$}
		{$\displaystyle\int F(x) \mathrm{\,d}x = f(x) + C$} 
		{\True $\displaystyle\int  F(x) \mathrm{\,d}x = x^2 - \ln|x| + C$}
	 \loigiai{
		Ta có $F(x) = \dfrac{2x^2-1}{x} = \dfrac{2x^2}{x} - \dfrac{1}{x} = 2x -  \dfrac{1}{x}$ (với $x \neq 0$).\\
		Vì $F(x)$ là  một nguyên hàm của $f(x)$, nên $f(x) = F'(x)=2 + \dfrac{1}{x^2}$.
		\begin{itemchoice}
			\itemch $f(x)  = 2 + \dfrac{1}{x^2}$.
			\itemch $\displaystyle\int f(x) \mathrm{\,d}x = \dfrac{2x^2-1}{x} +  C$.
			\itemch $\displaystyle\int  F(x) \mathrm{\,d}x = \displaystyle\int \left(2x - \dfrac{1}{x}\right) \mathrm{\,d}x = x^2 - \ln|x| + C$.
			\itemch $\displaystyle\int F(x) \mathrm{\,d}x = x^2 - \ln|x| + C$. 
		\end{itemchoice}
	} 
\end{ex}

% nguồn: Phan1_Lop_12_Dai_so_Chuong_4
\begin{ex}%[2D4H1-2]%[ma:64587173]
	Cho hàm số $F(x)=\displaystyle\int \left(\sqrt[5]{x^3}\right)\mathrm{\,d}x$  (với $x>0$). Khi đó
	 \choiceTF
		{$F(x)=\dfrac{3}{5}\displaystyle\int x\mathrm{\,d}x$} 
		{\True $F(x)+C=\displaystyle\int \left(\sqrt[5]{x^3}\right)\mathrm{\,d}x$, $C\in \mathbb{R}$} 
		 {$F(x)=\dfrac{3}{5}x^{-\tfrac{2}{5}}+C$}
		{Biết $F(1)=-\dfrac{3}{8}$, khi đó $F(x)=\dfrac{5}{8}\sqrt[5]{x^8}-\dfrac{3}{8}$} 
	 \loigiai{
		\begin{itemchoice}
			\itemch Ta có $F(x)=\displaystyle\int \left(\sqrt[5]{x^3}\right)\mathrm{\,d}x =  \displaystyle\int x^{\tfrac{3}{5}}\mathrm{\,d}x$.
			\itemch Nếu $F(x),G(x)$ là nguyên hàm của hàm  số $f(x)$ thì $F(x)=G(x)+C$.
			\itemch Ta có $F(x)=\displaystyle\int x^{\tfrac{3}{5}}\mathrm{\,d}x  = \dfrac{5}{8}x^{\tfrac{8}{5}}+C$.
			\itemch Ta có $F(x)=\dfrac{5}{8}x^{\tfrac{8}{5}}+C$, $F(1)=\dfrac{5}{8}+C=-\dfrac{3}{8} \Rightarrow C=-1  \Rightarrow F(x)=\dfrac{5}{8}\sqrt[5]{x^8}-1$.
		\end{itemchoice}
	} 
\end{ex}

% nguồn: Phan1_Lop_12_Dai_so_Chuong_4
\begin{ex}%[2D4H1-2]%[ma:bf4910c7]
	Cho hàm số $f(x)=\dfrac{2x+1}{x-1}$ với $x \neq  1$.
	 \choiceTF
		{\True  $f(x)=2+\dfrac{3}{x-1}$}
		{$\displaystyle\int  f(x)\mathrm{~d} x=2x+3\ln(x-1)+C$}
		{\True Nguyên hàm $F(x)$ của $f(x)=\dfrac{2x+1}{x-1}$ thỏa mãn  $F(2)=1$ là $F(x)=2x+3\ln|x-1|-3$}
		{\True Phương trình $F(x)=2x+2$ có $2$ nghiệm $x_1$; $x_2$. Khi đó  $T=x_1+x_2=2$}
	\loigiai{ 
		\begin{itemchoice}
			\itemch  Ta có $f(x)=\dfrac{2(x-1)+3}{x-1}=2+\dfrac{3}{x-1}$.
			\itemch $\displaystyle\int f(x)\mathrm{~d} x=\displaystyle\int\left(2+\dfrac{3}{x-1}\right)\mathrm{~d} x=2x+3\ln|x-1|+C$. 
			 \itemch Ta có $\displaystyle\int f(x)\mathrm{~d} x=\displaystyle\int\left(2+\dfrac{3}{x-1}\right)\mathrm{~d} x=2x+3\ln|x-1|+C$.\\
			Ta có $F(2)=1\Leftrightarrow  2\cdot 2+3\ln|2-1|+C=1\Leftrightarrow C=-3$ nên $F(x)=2x+3\ln|x-1|-3$.
			\itemch Ta có 
			\begin{eqnarray*}
				& &  F(x)=2x+2\\
				&\Leftrightarrow &2x+3\ln|x-1|-3=2x+2\\ 
				 &\Leftrightarrow &\ln|x-1|=5\Leftrightarrow|x-1|=\mathrm{e}^5\\
				&\Leftrightarrow  &
			\end{eqnarray*}$\Leftrightarrow \Leftrightarrow $\\
			$\Leftrightarrow\hoac{&x=1+\mathrm{e}^5\\ 
			 &x=1-\mathrm{e}^5}$
			(thỏa mãn điều kiện xác  định).\\
			 Vậy $T=x_1+x_2=\left(1+\mathrm{e}^5\right)+\left(1-\mathrm{e}^5\right)=2$.
		\end{itemchoice}
	 }
\end{ex}

% nguồn: Phan1_Lop_12_Dai_so_Chuong_4
\begin{ex}%[2D4V1-2]%[ma:7ba4671f]
	Hàm số $f(x)$ xác định trên $\mathbb{R}\setminus \{0\}$ thỏa mãn $f(x)=\dfrac{3x^5-2x^3+x^2+5x+4}{x^3}$. Khi đó 
	 \choiceTF
		{\True $f(x)=3x^2-2+\dfrac{1}{x}+\dfrac{1}{x^2}+\dfrac{4}{x^3}$} 
		{\True  $\displaystyle\int f(x)\mathrm{\,d}x=x^3-2x+\ln |x|-\dfrac{1}{x}-\dfrac{1}{2x^2}+C$}
		{\True Gọi $F(x)$ là một nguyên hàm của hàm số  $f(x)$ và thỏa mãn $F(-1)=0$. Khi đó tìm được $F(x)=x^3-2x+\ln |x|-\dfrac{1}{x}-\dfrac{1}{2x^2}-\dfrac{3}{2}$}
		{Gọi $G(x)$ là một nguyên hàm của hàm số $f(x)$. Biết $G(1)=0$ và $G(-3)-G(3)=2$. Khi đó tìm được $G(-2)=\dfrac{a}{b}+\dfrac{1}{\log_ce}$, với $a, b, c$ là các số nguyên dương và $\dfrac{a}{b}$ là  phân số tối giản. Vậy $a+b+c=1\,039$}
	\loigiai{ 
		\begin{itemchoice}
			\itemch Ta có  $f(x)=\dfrac{3x^5-2x^3+x^2+5x+4}{x^3}=3x^2-2+\dfrac{1}{x}+\dfrac{5}{x^2}+\dfrac{4}{x^3}$.
			\itemch Ta có  $\displaystyle\int f(x)\mathrm{\,d}x = \displaystyle\int\left(3x^2 - 2 + \dfrac{1}{x} + \dfrac{5}{x^2} + \dfrac{4}{x^3}\right)\mathrm{\,d}x = x^3 - 2x + \ln |x| - \dfrac{1}{x} - \dfrac{1}{2x^2} + C$.
			\itemch Với $F(-1)=0$ suy ra $C = -\dfrac{3}{2}$, nên $F(x) = x^3 - 2x + \ln |x| - \dfrac{1}{x} - \dfrac{1}{2x^2} -  \dfrac{3}{2}$.
			\itemch Ta có 
			\begin{itemize}
				\item Với $x>0$: $G(x) = x^3 - 2x  + \ln x - \dfrac{1}{x} - \dfrac{1}{2x^2} + C_1$, dùng $G(1)=0$ suy ra $C_1 = \dfrac{5}{2}$.
				\item Với $x<0$: $G(x) = x^3 - 2x + \ln(-x) - \dfrac{1}{x} - \dfrac{1}{2x^2}  + C_2$.
				\item Tính $G(-3)-G(3)=2$ suy  ra $C_2 = \dfrac{275}{6}$.
				\item Tính $G(-2) = -8 + 4 + \ln 2 + \dfrac{1}{2} -  \dfrac{1}{8} + \dfrac{275}{6} = \dfrac{1\,013}{24} + \dfrac{1}{\log_2 \mathrm{e}}$.
			\end{itemize}
			Vậy  $a=1\,013$, $b=24$, $c=2$ $\Rightarrow a + b + c = 1\,039$.
		\end{itemchoice}
	} 
\end{ex}

% nguồn: Phan1_Lop_12_Dai_so_Chuong_4
\begin{ex}%[2D4V1-2]%[ma:5bbd0df0]
	Hàm số $f(x)$  xác định trên $\mathbb{R}\setminus \{0\}$ thỏa mãn $f(x)=\dfrac{(x+1)(3x-2)}{x}$. Khi đó
	\choiceTF 
		{$f(x)=3x+5-\dfrac{2}{x}$} 
		{$\displaystyle\int f(x)\mathrm{\,d}x=3+x-2\ln  |x|+C$}
		{Gọi $F(x)$ là một nguyên hàm của hàm số $f(x)$ và thỏa mãn $F(3)=15$. Khi đó ta tìm được $F(x)=\dfrac{3x^2}{2}+x-2\ln |x|+15$} 
		{\True  Gọi $G(x)$ là một nguyên hàm của hàm số $f(x)$. Biết $G(2)=1$ và $G(5)+G(-5)=10$. Khi đó tìm được $G(-10)=a\ln 2+b\ln 5+c$, với $a, b, c$ là các số hữu tỷ. Vậy $a+b+c=75$}
	 \loigiai{
		\begin{itemchoice}
			\itemch  Ta có $f(x)=\dfrac{(x+1)(3x-2)}{x}=\dfrac{3x^2+x-2}{x}=3x+1-\dfrac{2}{x}$.
			 \itemch Ta có $\displaystyle\int f(x)\mathrm{\,d}x=\displaystyle\int \left(3x+1-\dfrac{2}{x}\right)\mathrm{\,d}x=\dfrac{3x^2}{2}+x-2\ln |x|+C$.
			\itemch Ta có $\displaystyle\int  f(x)\mathrm{\,d}x=\dfrac{3x^2}{2}+x-2\ln |x|+C$. \\
			 Do $F(3)=15 \Leftrightarrow \dfrac{27}{2}+3-2\ln 3+C=15 \Leftrightarrow C=2\ln 3 - \dfrac{3}{2}$. \\
			Vậy $F(x)=\dfrac{3x^2}{2}+x-2\ln  |x|+2\ln 3 - \dfrac{3}{2}$.
			\itemch Ta  có $\displaystyle\int f(x)\mathrm{\,d}x=\dfrac{3x^2}{2}+x-2\ln |x|+C = \heva{&\dfrac{3x^2}{2}+x-2\ln x+C_1 \text{ khi } x>0 \\ &\dfrac{3x^2}{2}+x-2\ln(-x)+C_2 \text{ khi } x<0.}$ \\
			 Do
			\begin{eqnarray*}
				&&\heva{&G(2)=1  \\ &G(5)+G(-5)=10}\\
				 &\Leftrightarrow&\heva{&6+2+C_1=1 \\ &\dfrac{75}{2}+5-2\ln 5+C_1+\dfrac{75}{2}-5-2\ln 5+C_2=10}\\
				&\Leftrightarrow&\heva{&C_1=-7 \\ &C_2=4\ln  5 - 65.}\\
			\end{eqnarray*}
			Nên $G(-10)=150-10-2\ln 10+4\ln 5 - 65=75 - 2\ln 2 + 2\ln 5 =  a\ln 2 + b\ln 5 + c$. \\
			Vậy  $a+b+c=-2+2+75=75$.
		\end{itemchoice}
	} 
\end{ex}

% nguồn: Phan1_Lop_12_Dai_so_Chuong_4
\begin{ex}%[2D4V1-2]%[ma:eccf6be8]
	 Cho hàm số $F(x)=x^2+x-6$ là một nguyên hàm của $f(x)$. Khi đó
	 \choiceTF
		 {$f(x)=\dfrac{x^3}{3}+\dfrac{x^2}{2}-6x+C$}
		 {$f(1)+f(2)+\cdots+f(50)=2400$}
		{Hàm  số $G(x)$ cũng là một nguyên hàm của $f(x)$ và $G(1)=3$ thì $G(4)=24$}
		{Hàm số $H(x-1)$ cũng là một nguyên hàm của  $f(x-1)$ và $H(0)=3$ thì $H(2)-H(4)=6$}
	 \loigiai{
		\begin{itemchoice}
			\itemch Ta có  $F'(x) = 2x+1$, do đó $f(x) = 2x+1$ khác với biểu thức đã cho.
			\itemch Ta  có $f(x) = 2x+1$ nên
			\begin{eqnarray*}
				 &f(1)+f(2)+\cdots+f(50)&= (2\cdot1+1)+(2\cdot2+1)+\cdots+(2\cdot50+1)\\
				&&= \sum\limits_{k=1}^{50} (2k+1)\\ 
				&&= 2\sum\limits_{k=1}^{50}k + \sum\limits_{k=1}^{50}1\\ 
				&&= 2\cdot\dfrac{50\cdot51}{2} + 50 \\ 
				&&= 2\,550  + 50 = 2\,600.
			\end{eqnarray*}
			\itemch Ta có $G(x) = x^2 + x + C$, $G(1)=3 \Rightarrow  1 + 1 + C = 3 \Rightarrow C = 1$. Vậy $G(4)=4^2 + 4 + 1 = 21$.
			\itemch Ta có $H(x-1) = (x-1)^2 + (x-1) + C =  x^2 - 2x + 1 + x - 1 + C = x^2 - x + C$. \\
			Biết $H(0)=3 \Rightarrow C = 3$ nên $H(x-1)  = x^2 - x + 3$. \\
			và $H(2) = H(3-1) = 3^2 - 3 + 3 = 9$, $H(4) = H(5-1) = 5^2 - 5 + 3 =  23$. \\
			Vậy $H(2)-H(4) = 9  - 23 = -14$.
		\end{itemchoice}
	} 
\end{ex}

% nguồn: Phan1_Lop_12_Dai_so_Chuong_4
\begin{ex}%[2D4V1-2]%[ma:fffc13d1]
	Hàm hai số $f(x)$ và  $g(x)$ xác định trên $\mathbb{R} \setminus \{0\}$ thỏa mãn $f(x)=\dfrac{x^2+1}{x}$ và $\displaystyle\int g(x)\mathrm{\,d}x = x + \ln |x| + C$. Khi đó
	 \choiceTF
		{\True $f(x) = x  + \dfrac{1}{x}$}
		{$g(x) = 1 + \dfrac{1}{x} +  C$}
		{$F(x)$ là một nguyên  hàm của $f(x)+g(x)$ và thỏa mãn $F(-1)=3$. Khi đó tìm được $F(x)=x^2+2\ln|x|-\dfrac{1}{2}$}
		{\True $G(x)$ là  một nguyên hàm của $f(x)+g(x)$. Biết $G(2)=2\ln 2$ và $G(4)+G(-4)=2$. Khi đó tìm được $G(-6)=a\ln 2 + b\ln 3 + c$, với $a, b, c$ là các số thực. Vậy $a+b+c=-2$}
	\loigiai{ 
		\begin{itemchoice}
			\itemch Ta có $f(x) = \dfrac{x^2+1}{x} = x + \dfrac{1}{x}$. 
			 \itemch Ta có $\displaystyle\int g(x)\mathrm{\,d}x = x + \ln|x| + C \Rightarrow g(x) = 1 + \dfrac{1}{x}$ (không có $+C$).
			\itemch Ta có  $f(x)+g(x) = x + \dfrac{1}{x} + 1 + \dfrac{1}{x} = x + 1 + \dfrac{2}{x}$.\\
			Do đó $\displaystyle\int  (f+g)\mathrm{\,d}x = \dfrac{x^2}{2} + x + 2\ln|x| + C$.\\
			 Vì $F(-1)=3$ nên $C = \dfrac{7}{2}$. Do đó $F(x) = \dfrac{x^2}{2} + x + 2\ln|x| + \dfrac{7}{2}$.
			 \itemch
			\begin{itemize}
				\item Với  $x > 0$: $G(x) = \dfrac{x^2}{2} + x + 2\ln x + C_1$, vì $G(2) = 2\ln 2$ $\Rightarrow C_1 = -4$.
				\item Với $x < 0$: $G(x) = \dfrac{x^2}{2} + x +  2\ln(-x) + C_2$.
				\item Vì  $G(4)+G(-4)=2 \Rightarrow C_2 = -8\ln 2 -10$.
				\item Ta có $G(-6) = \dfrac{36}{2} -6 + 2\ln 6 + C_2 = 18 - 6 + 2\ln 6 -8\ln  2 -10 = 2 - 6\ln 2 + 2\ln 3$.
			\end{itemize}
			Do đó $a = -6$, $b = 2$, $c = 2 \Rightarrow  a + b + c = -2$.
		\end{itemchoice}
	} 
\end{ex}

% nguồn: Phan1_Lop_12_Dai_so_Chuong_4
\begin{ex}%[2D4V1-2]%[ma:0e9ab05b]
	Hàm số $f(x)$  xác định trên $\mathbb{R}\setminus \{0\}$ thỏa mãn $f(x)=\dfrac{x^2+5x-7}{x}$. Khi đó
	 \choiceTF
		{\True $f(x)=x+5-\dfrac{7}{x}$} 
		{\True $\displaystyle\int f(x)\mathrm{\,d}x=\dfrac{x^2}{2}+5x-7\ln |x|+C$} 
		{Gọi $F(x)$ là một nguyên hàm của hàm số $f(x)$ và thỏa mãn $F(1)=5$. Khi đó ta tìm được hàm số $F(x)=\dfrac{x^2}{2}+5x-7\ln  |x|+\dfrac{1}{2}$}
		{Gọi $G(x)$ là một nguyên hàm của hàm số $f(x)$. Biết $G(1)=4$ và $G(3)+G(-9)=20$. Khi đó tìm được $G(-6)=a\ln 2+b\ln 3+c$, với $a, b, c$ là các  số hữu tỉ. Vậy $a+b+c=\dfrac{2}{3}$}
	 \loigiai{
		\begin{itemchoice}
			\itemch  Ta có $f(x)=\dfrac{x^2+5x-7}{x}=x+5-\dfrac{7}{x}$.
			\itemch  Ta có $\displaystyle\int f(x)\mathrm{\,d}x=\displaystyle\int \left(x+5-\dfrac{7}{x}\right)\mathrm{\,d}x=\dfrac{x^2}{2}+5x-7\ln |x|+C$.
			\itemch Ta có $\displaystyle\int f(x)\mathrm{\,d}x=\dfrac{x^2}{2}+5x-7\ln  |x|+C$ \\
			Do $F(1)=5\Leftrightarrow \dfrac{1}{2}+5+C=5\Leftrightarrow C=-\dfrac{1}{2}$  nên $F(x)=\dfrac{x^2}{2}+5x-7\ln |x|-\dfrac{1}{2}$.
			\itemch  Ta có $\displaystyle\int f(x)\mathrm{\,d}x=\dfrac{x^2}{2}+5x-7\ln |x|+C=\heva{&\dfrac{x^2}{2}+5x-7\ln x+C_1 \text{ khi } x>0 \\ &\dfrac{x^2}{2}+5x-7\ln(-x)+C_2 \text{ khi } x<0.}$ \\
			Do 
			\begin{eqnarray*}
				&&\heva{&G(1)=4 \\ &G(3)+G(-9)=20}\\ 
				&\Leftrightarrow&\heva{&C_1=-\dfrac{3}{2} \\ &\dfrac{9}{2}+15-7\ln  3-\dfrac{3}{2}+\dfrac{81}{2}-45-14\ln 3+C_2=20}\\
				&\Leftrightarrow&\heva{&C_1=-\dfrac{3}{2} \\  &C_2=21\ln 3+\dfrac{13}{2}.}
			\end{eqnarray*}
			Khi đó  $G(-6)=\dfrac{36}{2}-30-7\ln 6+21\ln 3+\dfrac{13}{2}=-7\ln 2+14\ln 3-\dfrac{11}{2}$. \\
			Vậy $a+b+c=-7+14-\dfrac{11}{2}=\dfrac{3}{2}$. 
		\end{itemchoice}
	} 
\end{ex}

% nguồn: Phan1_Lop_12_Dai_so_Chuong_4
\begin{ex}%[2D4V1-2]%[ma:ca7162d1]
	Cho hàm số $f(x)=3$. Khi  đó
	\choiceTF 
		{\True $\displaystyle\int  f(x)\mathrm{\,d}x = 3x+C$}
		{$\displaystyle\int \left[f(x)+x\right]^2\mathrm{\,d}x  = x^3-\dfrac{1}{2}x^2+2x+C$}
		{Gọi $F(x)$ là một nguyên hàm của  $f(x)$. Biết $F(1)=1$ thì $F(x)=3x-1$}
		 {Gọi $F(x)$ là một nguyên hàm của $f(x)$ thì $F(1)+F(2)+\cdots+F(100)=14\,590$ khi $F(1)=1$}
	\loigiai{ 
		\begin{itemchoice}
			 \itemch Ta có $\displaystyle\int 3\mathrm{\,d}x = 3x+C$.
			\itemch Ta có $\displaystyle\int \left[f(x)+x\right]^2\mathrm{\,d}x = \displaystyle\int (3+x)^2\mathrm{\,d}x =  \displaystyle\int (x^2+6x+9)\mathrm{\,d}x = \dfrac{x^3}{3}+3x^2+9x+C$.
			\itemch Ta có $F(1)=1 \Leftrightarrow 3+C=1 \Leftrightarrow C=-2$  nên $F(x)=3x-2$.
			\itemch Ta có $F(x)=3x-2$ nên $F(1)=1$, $F(2)=4$, $F(3)=7$, \ldots, $F(100)=298$ là cấp số  cộng với $u_1=1$, $d=3$. \\
			Khi  đó $F(1)+F(2)+\cdots+F(100)=\dfrac{(1+298)\cdot 100}{2}=14\,950$.
		\end{itemchoice}
	} 
\end{ex}

% nguồn: Phan1_Lop_12_Dai_so_Chuong_4
\begin{ex}%[2D4V1-2]%[ma:f0558875]
	Cho hàm số $f(x)=x+1$.  Khi đó
	 \choiceTF
		{\True $\displaystyle\int f(x)\mathrm{\,d}x =  x^2+x+C$}
		{\True $\displaystyle\int \left[(x-1)f(x)\right]\mathrm{\,d}x = \dfrac{1}{3}x^3 - x + C$} 
		{Gọi $F(x)$ là một nguyên hàm của $f(x)$. Biết $F(1)=2$ thì $F(x)=x^2+x-1$} 
		{\True Gọi $F(x)$ là một nguyên  hàm của $f(x)$. Biết $F(1)=2$ và $\dfrac{1}{F(1)}+\dfrac{1}{F(2)}+\cdots+\dfrac{1}{F(100)} = \dfrac{a}{b}$ thì $a+b=201$}
	 \loigiai{
		\begin{itemchoice}
			\itemch Ta có $\displaystyle\int  f(x)\mathrm{\,d}x = \displaystyle\int (x+1)\mathrm{\,d}x = x^2+x+C$.
			\itemch  Ta có $\displaystyle\int \left[(x-1)f(x)\right]\mathrm{\,d}x = \displaystyle\int (x-1)(x+1)\mathrm{\,d}x = \displaystyle\int (x^2-1)\mathrm{\,d}x = \dfrac{1}{3}x^3 - x + C$.
			\itemch Ta có $F(x) = x^2 + x + C$, $F(1)=2 \Rightarrow 1^2 + 1 + C = 2 \Rightarrow C =  0$ nên $F(x) = x^2 + x$.
			\itemch Ta có $F(x) = x^2 + x =  x(x+1)$ nên $\dfrac{1}{F(x)}= \dfrac{1}{x(x+1)}= \dfrac{1}{x} - \dfrac{1}{x+1}$. Do đó
			\begin{eqnarray*}
				&\dfrac{1}{F(1)}+\dfrac{1}{F(2)}+\cdots+\dfrac{1}{F(100)}&= \left(1 - \dfrac{1}{2}\right) + \left(\dfrac{1}{2} - \dfrac{1}{3}\right) + \cdots  + \left(\dfrac{1}{100} - \dfrac{1}{101}\right)\\
				&&= 1 - \dfrac{1}{101}  = \dfrac{100}{101}.
			\end{eqnarray*}
			Suy  ra $a=100$, $b=101 \Rightarrow a+b=201$.
		\end{itemchoice}
	} 
\end{ex}

% nguồn: Phan1_Lop_12_Dai_so_Chuong_4
\begin{ex}%[2D4V1-2]%[ma:034d7982]
	Hàm số $f(x)$ xác định  trên $\mathbb{R}\setminus \{0\}$ thỏa mãn $f(x)=\left(\dfrac{x-1}{x}\right)^2$. Khi đó
	\choiceTF 
		 {\True $f(x)=1-\dfrac{2}{x}+\dfrac{1}{x^2}$}
		{$\displaystyle\int f(x)\mathrm{\,d}x = x-2\ln |x|+\dfrac{1}{x}+C$} 
		{$F(x)$ là một nguyên hàm  của $f(x)$ và $F(3)=5$. Khi đó $F(x)=x-2\ln|x|-\dfrac{1}{x}+5$}
		{\True $G(x)$ là một nguyên hàm của $f(x)$, biết $G(1)=5$ và $G(2)+G(-4)=2025$. Khi đó $G(-8)=\dfrac{a}{b}$  với $\dfrac{a}{b}$ tối giản. Vậy $a+b=16\,123$}
	\loigiai{ 
		\begin{itemchoice}
			\itemch Ta có  $f(x)=\left(\dfrac{x-1}{x}\right)^2 = \dfrac{x^2 - 2x + 1}{x^2} = 1 - \dfrac{2}{x} + \dfrac{1}{x^2}$.
			\itemch Ta có $\displaystyle\int f(x)\mathrm{\,d}x = \displaystyle\int \left(1 - \dfrac{2}{x} + \dfrac{1}{x^2}\right)\mathrm{\,d}x = x - 2\ln|x|  - \dfrac{1}{x} + C$.
			\itemch  Vì $F(3)=5 \Rightarrow 3 - 2\ln 3 - \dfrac{1}{3} + C = 5 \Rightarrow C = 2\ln 3 + \dfrac{7}{3}$.\\
			Vậy $F(x) = x -  2\ln|x| - \dfrac{1}{x} + 2\ln 3 + \dfrac{7}{3}$.
			 \itemch
			\begin{itemize}
				\item Với $x > 0$: $G(x) = x - 2\ln x - \dfrac{1}{x} + C_1$, từ $G(1)=5$ suy ra $C_1 = 5$. 
				\item Với  $x < 0$: $G(x) = x - 2\ln(-x) - \dfrac{1}{x} + C_2$.
				\item Ta  có
				\begin{align*}
					& G(2) = 2 - 2\ln 2 - \dfrac{1}{2} + 5 = \dfrac{9}{2}  - 2\ln 2 \\
					& G(-4) = -4 - 2\ln 4 + \dfrac{1}{4}  + C_2 = -4 - 4\ln 2 + \dfrac{1}{4} + C_2.
				\end{align*}
				Ta có $G(2) + G(-4) =  2\,025$ nên
				\begin{align*}
					G(2) + G(-4)& = \dfrac{9}{2} - 2\ln  2 -4 - 4\ln 2 + \dfrac{1}{4} + C_2\\
					& = \dfrac{7}{4} -6\ln 2 + C_2 = 2\,025  \\
					& \Rightarrow  C_2 = 6\ln 2 + \dfrac{8\,089}{4}.
				\end{align*}
				\item Tính  $G(-8) = -8 - 6\ln 2 + \dfrac{1}{8} + 6\ln 2 + \dfrac{8\,089}{4} = \dfrac{16\,115}{8}$.
			\end{itemize}
			Suy ra $a = 16\,115$, $b = 8 \Rightarrow a + b = 16\,123$. 
		\end{itemchoice}
	} 
\end{ex}


%%%%%%%%%%%%%%%% TRẢ LỜI NGẮN %%%%%%%%%%%%%%%%

% nguồn: Phan2_Lop_12_Dai_so_Chuong_4
\begin{bt}%[2D4H1-2]%[ma:6dabb024]
	Cho hàm số  $f(x)=x^2+4x$. Biết $\displaystyle\int f(x)\mathrm{\,d}x=ax^3+bx^2+C$, trong đó $a$, $b\in\mathbb{R}$. Tính $3a+b$.
	 \shortans{3}
	\loigiai{ 
		 Ta có $\displaystyle\int f(x)\mathrm{\,d}x=\displaystyle\int(x^2+4x)\mathrm{\,d}x=\dfrac{1}{3}x^3+2x^2+C$.\\
		Suy  ra $a=\dfrac{1}{3}$, $b=2$ nên $3a+b=3\cdot\dfrac{1}{3}+2=3$.
	} 
\end{bt}

% nguồn: Phan2_Lop_12_Dai_so_Chuong_4
\begin{bt}%[2D4H1-2]%[ma:58bdb86b]
	Một nguyên hàm của hàm $f(x)=5x^2-2x\sqrt{x}-\dfrac{1}{x}$ có  dạng $F(x)=ax^3+bx^{\frac{5}{2}}-\ln x+2\,025$. Tính giá trị biểu thức $P=6ab$.
	\shortans{-8} 
	\loigiai{ 
		Ta có $$\displaystyle\int{f(x)}\,\mathrm{\,d}{x} 
		 =\displaystyle\int{\left(5x^2-2x^{\frac{3}{2}}-\dfrac{1}{x}\right)}\mathrm{\,d}{x}
		=\dfrac{5}{3}x^3-\dfrac{4}{5}x^{\frac{5}{2}}-\ln x+C.$$ 
		Như vậy một nguyên hàm của $f(x)$ thỏa mãn yêu cầu  đề bài là
		$$F(x)=\dfrac{5}{3}x^3-\dfrac{4}{5}x^{\frac{5}{2}}-\ln x+2\,025.$$ 
		 Do đó $a=\dfrac{5}{3}$ và $b=-\dfrac{4}{5}$. \\
		Vậy $6ab=6  \cdot \dfrac{5}{3} \cdot \left(-\dfrac{4}{5}\right)=-8$.
	 }
\end{bt}

% nguồn: Phan2_Lop_12_Dai_so_Chuong_4
\begin{bt}%[2D4H1-2]%[ma:8736141c]
	Cho hàm số $f(x)=3x^2+2x$. Biết rằng hàm số $f(x)$ có một nguyên hàm là hàm số $F(x)=ax^3+bx^2+c$.  Tính giá trị của biểu thức $T=a^2+b^2$.
	 \shortans{2}
	 \loigiai{
		Ta có $F(x) = \displaystyle\int f(x) \mathrm{\,d}x = \int (3x^2 + 2x) \mathrm{\,d}x = x^3 +  x^2 + c$. Vậy $a=1$, $b=1$ $\Rightarrow T=2$.
	 }
\end{bt}

% nguồn: Phan2_Lop_12_Dai_so_Chuong_4
\begin{bt}%[2D4H1-2]%[ma:4960f690]
	Hàm số  $f(x)=\dfrac{1}{x^2(x+1)}$ có một nguyên hàm là $F(x)$ thỏa $F(1)=\ln 2$. Tính $F(-2)$ (Kết quả làm tròn đến phần chục).\\
	 \shortans{0,8}
	\loigiai{ 
		Ta có 
		\begin{eqnarray*}
			 &f(x)&=\dfrac{1}{x^2(x+1)}=\dfrac{1+x-x}{x^2(x+1)} =\dfrac{1}{x^2}-\dfrac{1}{x(x+1)}\\
			 &&=\dfrac{1}{x^2}-\left(\dfrac{1}{x}-\dfrac{1}{x+1}\right)=\dfrac{1}{x^2}-\dfrac{1}{x}+\dfrac{1}{x+1}.
		\end{eqnarray*}
		 Do đó,
		\begin{eqnarray*}
			&F(x)&=\displaystyle\int  f(x)\mathrm{\,d}x=\displaystyle\int\left(\dfrac{1}{x^2}-\dfrac{1}{x}+\dfrac{1}{x+1}\right)\mathrm{\,d}x\\
			&&=-\dfrac{1}{x}-\ln  |x|+\ln |x+1|+C=-\dfrac{1}{x}+ \ln\left|\dfrac{x+1}{x}\right|+C.
		\end{eqnarray*}
		Ta có $F(1)=-\dfrac{1}{1}+\ln\left|\dfrac{1+1}{1}\right|  +C=-1+\ln 2+C=\ln 2$.\\
		 Suy ra $C = 1$.\\
		Vậy $F(x) = -\dfrac{1}{x} + \ln\left| \dfrac{x+1}{x} \right|  + 1 $.\\
		$F(-2)=-\dfrac{1}{-2}+\ln\left|\dfrac{-2+1}{-2}\right|+1=\dfrac{1}{2}+\ln\left|\dfrac{-1}{-2}\right|+1=\dfrac{3}{2}+\ln\dfrac{1}{2}=\dfrac{3}{2}-\ln 2$.\\ 
		Ta có  $\ln 2 \approx 0{,}693$.\\
		$F(-2)=1{,}5-0{,}693=0{,}807\approx 0{,}8$. 
	 }
\end{bt}

% nguồn: Phan2_Lop_12_Dai_so_Chuong_4
\begin{bt}%[2D4H1-2]%[ma:44823784]
	Cho hàm số $f(x)=\dfrac{1}{x}$ có nguyên hàm  là $F(x)$ thỏa mãn $F(1)=1$. Tính $F(\sqrt{\mathrm{e}})$.
	 \shortans{1,5}
	\loigiai{ 
		Ta có $F(x) = \displaystyle\int f(x)\mathrm{\, d}x = \displaystyle\int \dfrac{1}{x}\mathrm{\, d}x  = \ln|x| + C$.\\
		Vì $F(1)=1$ nên $\ln|1|+C = 1 \Rightarrow 0+C =  1 \Rightarrow C = 1$.
		Do đó  $F(x) = \ln|x|+1$.\\
		Vậy $F(\sqrt{\mathrm{e}}) = \ln|\sqrt{\mathrm{e}}|+1 = \ln(\mathrm{e}^{\frac{1}{2}})+1 = \dfrac{1}{2}\ln\mathrm{e}+1 = \dfrac{1}{2}  \cdot 1+1 = 0{,}5+1 = 1{,}5$.
	} 
\end{bt}

% nguồn: Phan2_Lop_12_Dai_so_Chuong_4
\begin{bt}%[2D4H1-2]%[ma:e7b23c72]
	Cho hàm số $F(x)$ là một nguyên hàm của hàm số $f(x)=\left(x^2-2\right)(2x+1)$ và $F(-1)=\dfrac{1}{6}$. Tính $F\left(-\dfrac{1}{2}\right)$ (viết kết quả dưới dạng số thập phân và làm tròn  đến hàng phần~trăm).
	 \shortans{0,49}
	 \loigiai{
		Ta có $f(x)=\left(x^2-2\right)(2  x+1)=2 x^3+x^2-4 x-2$.\\
		Suy ra 
		 \[
		\begin{aligned}
			F(x)=\displaystyle\int f(x) \mathrm{\,d} x&=\displaystyle\int\left(2  x^3+x^2-4x-2\right) \mathrm{\,d} x\\
			 & =\displaystyle\int 2x^3 \mathrm{\,d} x+\displaystyle\int x^2 \mathrm{\,d} x-\displaystyle\int 4x \mathrm{\,d}x-\displaystyle\int 2 \mathrm{\,d} x \\
			 & =\dfrac{1}{2} x^4+\dfrac{1}{3} x^3-2 x^2-2 x+C,\, C \in \mathbb{R}.
		\end{aligned}
		\] 
		Mà  $F(-1)=\dfrac{1}{6}$ nên suy ra $C=0$. \\
		Vậy hàm số $F(x)=\dfrac{1}{2} x^4+\dfrac{1}{3} x^3-2 x^2-2  x$.\\
		Suy ra $F\left(-\dfrac{1}{2}\right)=\dfrac{47}{96}  \approx 0{,}49$.\\
		\textbf{Cách khác:}\\ 
		$I=\displaystyle\int_{-1}^{-\tfrac{1}{2}} f(x) \mathrm{\,d}x=F\left(-\dfrac{1}{2}\right)-F(-1)\Rightarrow F\left(-\dfrac{1}{2}\right)=I+F(-1)=\dfrac{47}{96}\approx  0{,}49.$
	} 
\end{bt}


%%%%%%%%%%%%%%%% TỰ LUẬN %%%%%%%%%%%%%%%%

% nguồn: Phan2_Lop_12_Dai_so_Chuong_4
\begin{bt}%[2D4N1-2]%[ma:52258946]
	Tìm họ nguyên hàm của hàm  số $f(x)$, biết
	\begin{listEX}[2]
		 \item $f(x)=x^2$
		 \item $f(x)=x^4+x^{-3}$
		\item  $f(x)=x^2+\dfrac{3}{x}$
		\item $f(x)=\dfrac{1}{x}-\dfrac{1}{x^2}$ 
		\item  $f(x)=x^2+\dfrac{3}{x}-2\sqrt{x}$
		\item $f(x)=3\sqrt{x}+\dfrac{1}{\sqrt[3]x}$ 
	\end{listEX}
	\loigiai{ 
		\begin{enumerate}
			\item Ta có  $\displaystyle\int x^2\mathrm{d} x=\dfrac{x^3}{3}+C$.
			\item Ta có $\displaystyle\int f(x)\mathrm{\,d}x=\displaystyle\int\left(x^4+x^{-3}\right)\mathrm{\,d}x=\dfrac{1}{5} x^5-\dfrac{1}{2} x^{-2}+C$. 
			 \item Ta có $\displaystyle\int\left(x^2+\dfrac{3}{x}\right)\mathrm{\,d}x=\dfrac{x^3}{3}+3\ln|x|+C$.
			 \item Ta có $\displaystyle\int\left(\dfrac{1}{x}-\dfrac{1}{x^2}\right)\mathrm{\,d}x=\ln|x|+\dfrac{1}{x}+C$.
			\item Ta có $\displaystyle\int x^2+\dfrac{3}{x}-2\sqrt{x}  \mathrm{d} x=\dfrac{x^3}{3}+3\ln|x|-\dfrac{4}{3} \sqrt{x^3}+C$.
			\item  Ta có $\displaystyle\int\left(3\sqrt{x}+\dfrac{1}{\sqrt[3]x}\right)\mathrm{\,d}x=\displaystyle\int\left(3x^{\tfrac{1}{2}}+x^{-\tfrac{1}{3}}\right)\mathrm{\,d}x=\dfrac{6}{3} x^{\tfrac{3}{2}}+\dfrac{3}{2} x^{\tfrac{2}{3}}+C=2x^{\tfrac{3}{2}}+\dfrac{3}{2} x^{\tfrac{2}{3}}+C$.
		\end{enumerate}
	 }
\end{bt}

% nguồn: Phan2_Lop_12_Dai_so_Chuong_4
\begin{bt}%[2D4H1-2]%[ma:77fb388d]
	Tìm  họ nguyên hàm của hàm số $f(x)$, biết
	\begin{listEX}[2]
		\item  $f(x)=\dfrac{x^2+3x+2}{x}$.
		\item $f(x)=\dfrac{x^2+3x+1}{\sqrt{x}}$. 
	\end{listEX}
	 \loigiai{
		\begin{enumerate}
			\item $\displaystyle\int f(x)\mathrm{\,d}x=\displaystyle\int \dfrac{x^2+3x+2}{x} \mathrm{d} x=\displaystyle\int\left(x+3+\dfrac{2}{x}\right)\mathrm{\,d}x=\dfrac{1}{2}  x^2+3x+2\ln|x|+C$.
			\item Ta có 
			\begin{eqnarray*}
				\displaystyle\int f(x)\mathrm{\,d}x&= &\displaystyle\int \dfrac{x^2+3x+1}{\sqrt{x}} \mathrm{d} x=\displaystyle\int\left(x  \sqrt{x}+3\sqrt{x}+\dfrac{1}{\sqrt{x}}\right)\mathrm{\,d}x\\
				 &=&\displaystyle\int\left(x^{\tfrac{3}{2}}+3x^{\tfrac{1}{2}}+\dfrac{1}{\sqrt{x}}\right)\mathrm{\,d}x\\
				&=& \dfrac{2}{5} x^{\tfrac{5}{2}}+2x^{\tfrac{3}{2}}+2\sqrt{x}+C. 
			\end{eqnarray*}
		\end{enumerate}
		Vậy  $\displaystyle\int f(x)\mathrm{\,d}x=\dfrac{2}{5} x^{\tfrac{5}{2}}+2x^{\tfrac{3}{2}}+2\sqrt{x}+C$.
	} 
\end{bt}

% nguồn: Phan2_Lop_12_Dai_so_Chuong_4
\begin{bt}%[2D4H1-2]%[ma:2f86f82d]
	Tìm nguyên hàm của hàm số $f(x)$,  biết
	\begin{listEX}[2]
		 \item $f(x)=2x\left(1+3x^3\right)$
		\item  $f(x)=(3x-2)^2$
		 \item $f(x)=\dfrac{2x^4+3}{x^2}$
		 \item $f(x)=\dfrac{1}{2x-1}$
		 \item $f(x)=\dfrac{2}{4x-3}$
		\item $f(x)=\dfrac{6x+2}{3x-1}$ 
	\end{listEX}
	\loigiai{ 
		\begin{enumerate}
			 \item Ta có $\displaystyle\int 2x\left(1+3x^3\right)\mathrm{\,d}x=\displaystyle\int\left(2x+6x^4\right)\mathrm{\,d}x=x^2\left(1+\dfrac{6x^3}{5}\right)+C$.
			\item Ta có  $\displaystyle\int(3x-2)^2\mathrm{d} x=\displaystyle\int\left(9x^2-12x+4\right)\mathrm{\,d}x=3x^3-6x^2+4x+C$.
			\item Ta có $f(x)=\dfrac{2x^4+3}{x^2}=2x^2+3x^{-2} \Rightarrow \displaystyle\int f(x)\mathrm{\,d}x=\displaystyle\int\left(2x^2+3x^{-2}\right)\mathrm{\,d}x=\dfrac{2x^3}{3}-\dfrac{3}{x}+C$. 
			 \item Ta có $F(x)=\displaystyle\int f(x)\mathrm{\,d}x=\displaystyle\int \dfrac{1}{2x-1} \mathrm{d} x=\dfrac{1}{2} \ln|2x-1|+C$.
			\item Ta có $\displaystyle\int \dfrac{2}{4x-3} \mathrm{d} x=\displaystyle\int \dfrac{1}{2x-\dfrac{3}{2}} \mathrm{d} x=\dfrac{1}{2}  \ln\left|2x-\dfrac{3}{2}\right|+C$.
			\item Ta có $\displaystyle\int \dfrac{6x+2}{3x-1} \mathrm{d} x=\displaystyle\int\left(2+\dfrac{4}{3x-1}\right)\mathrm{\,d}x=2x+\dfrac{4}{4} \ln|3x-1|+C$. 
		\end{enumerate}
	} 
\end{bt}

% nguồn: Phan2_Lop_12_Dai_so_Chuong_4
\begin{bt}%[2D4H1-2]%[ma:431ea150]
	 Tìm $\displaystyle\int\left( x^2-3\sqrt{x}\right) \mathrm{\,d}x$.
	\loigiai{ 
		Ta có $\displaystyle\int(x^2-3\sqrt{x})\mathrm{\,d}x=\displaystyle\int  x^2\mathrm{\,d}x-3\displaystyle\int x^{\frac{1}{2}}\mathrm{\,d}x=\dfrac{x^3}{3}-3\dfrac{x^{\frac{3}{2}}}{\dfrac{1}{2}+1}+C=\dfrac{x^3}{3}-2x^{\frac{3}{2}}+C.$
	 }
\end{bt}

% nguồn: Phan2_Lop_12_Dai_so_Chuong_4
\begin{bt}%[2D4H1-2]%[ma:ed43b920]
	Cho hàm số $f(x)$ có đạo hàm $f'(x)=\dfrac{3x^3+1}{x}$. Khi  đó hãy tính giá trị biểu thức $f(-3)+f(3)$ biết $f(1)=2$, $f(-1)=3$.
	\loigiai{ 
		\begin{eqnarray*}
			f(x)  &= &\displaystyle\int f'(x)\mathrm{\,d}x=\displaystyle\int \dfrac{3x^3+1}{x} \mathrm{d} x\\
			&=& \displaystyle\int\left(3x^2+\dfrac{1}{x}\right)\mathrm{\,d}x=x^3+\ln|x|+C=\heva{&x^3+\ln x+C_1\text{ khi } x > 0\\ 
			&x^3+\ln(-x)+C_2\text{ khi } x <  0}
		\end{eqnarray*}
		 $f(1)=1+C_1=2\Rightarrow C_1=1; f(-1)=-1+C_2=3\Rightarrow C_2=4$.\\
		Suy ra $f(x)=\heva{&x^3+\ln x+1\,\text{ khi  } x > 0\\
		&x^3+\ln(-x)+4\,\text{ khi } x  < 0}$ nên $f(-3)+f(3)=2\ln 3+5$.
	} 
\end{bt}

% nguồn: Phan2_Lop_12_Dai_so_Chuong_4
\begin{bt}%[2D4H1-2]%[ma:790ac3f4]
	Biết $F(x)$ là một  nguyên hàm của hàm số $f(x)=\dfrac{1+2x^2}{x}$ thỏa mãn $F(-1)=3$. Tìm $F(x)$.
	 \loigiai{
		Ta  có $F(x)=\displaystyle\int \dfrac{1+2x^2}{x} \mathrm{d} x=\displaystyle\int\left(\dfrac{1}{x}+2x\right)\mathrm{\,d}x=\ln|x|+x^2+C$.\\
		Mà $F(-1)=3\Leftrightarrow \ln|-1|+(-1)^2+C=3\Leftrightarrow C=2$ nên  $F(x)=\ln|x|+x^2+2$.
	 }
\end{bt}

% nguồn: Phan2_Lop_12_Dai_so_Chuong_4
\begin{bt}%[2D4H1-2]%[ma:7165020a]
	Biết $F(x)$ là một nguyên hàm của  hàm số $f(x)=\dfrac{1}{2x+1}$ và $F(0)=2$. Tính $F(\mathrm{e})$.
	\loigiai{ 
		Ta có  $F(x)=\displaystyle\int f(x)\mathrm{\,d}x=\dfrac{1}{2} \ln|2x+1|+C$.\\
		Do $F(0)=2$ suy  ra $C=2$ nên $F(\mathrm{e})=\dfrac{1}{2} \ln(2\mathrm{e}+1)+2=\ln\sqrt{2\mathrm{e}+1}+2$.
	 }
\end{bt}

% nguồn: Phan2_Lop_12_Dai_so_Chuong_4
\begin{bt}%[2D4H1-2]%[ma:a82863f9]
	Cho hàm số $y=f(x)$ có đạo hàm  là $f'(x)=6x^2-2, \forall x \in \mathbb{R}$ và $f(0)=1$. Biết $F(x)$ là một nguyên hàm của $f(x)$ và $F(2)=0$. Tính $F(4)$.
	 \loigiai{
		Ta có $\displaystyle\int  f'(x)\mathrm{\,d}x=\displaystyle\int\left(6x^2-2\right)\mathrm{\,d}x=2x^3-2x+C$.
		Suy ra  $f(x)=2x^3-2x+C$.\\
		Theo bài ra ta  có $f(0)=1\Rightarrow C=1\Rightarrow f(x)=2x^3-2x+1$.\\
		Khi đó $F(x)=\displaystyle\int f(x)\mathrm{\,d}x=\displaystyle\int\left(2x^3-2x+1\right)\mathrm{\,d}x=\dfrac{1}{2}  x^4-x^2+x+C'$.\\
		Mà $F(2)=0\Rightarrow 6+C'=0\Rightarrow C'=-6$ suy ra $F(x)=\dfrac{1}{2}  x^4-x^2+x-6\Rightarrow F(4)=110$.\\
		Vậy  $F(4)=110$.
	 }
\end{bt}

% nguồn: Phan2_Lop_12_Dai_so_Chuong_4
\begin{bt}%[2D4H1-2]%[ma:66ded271]
	Tìm giá trị  của tham số $m$ biết $F(x)$ là nguyên hàm của hàm số $\break f(x)=\dfrac{1}{2\sqrt{x+1}}+m-1$ thỏa mãn $F(0)=0$ và $F(3)=7$.
	 \loigiai{
		Ta  có
		\begin{eqnarray*}
			 F(x)& =& \displaystyle\int\left(\dfrac{1}{2\sqrt{x+1}}+m-1\right)\mathrm{\,d}x\\
			 &=&\displaystyle\int\left(\dfrac{1}{2\sqrt{x+1}}\right)\mathrm{\,d}x+\displaystyle\int(m-1)\mathrm{\,d}x\\
			&=  &\displaystyle\int\left(\dfrac{1}{2\sqrt{x+1}}\right) d(x+1)+\displaystyle\int(m-1)\mathrm{\,d}x=\sqrt{x+1}+(m-1) x+C.
		\end{eqnarray*}
		Theo  đề bài ta có $\heva{&F(0)=0\\
		&F(3)=7} \Rightarrow\heva{&C+1=0\\ 
		 &C+3m=8} \Leftrightarrow\heva{&C=-1\\
		&m=3}$.\\ 
		Vậy $m=3$ khi  đó $F(x)=\sqrt{x+1}+2x-1$.
	 }
\end{bt}

% nguồn: Phan2_Lop_12_Dai_so_Chuong_4
\begin{bt}%[2D4H1-2]%[ma:980e7965]
	 Cho hàm số $f(x)=\heva{&2x-1\text{ khi } x \geq 1\\
	&3x^2-2\text{ khi } x < 1.}$ Giả sử $F$ là nguyên hàm của $f$ trên R thỏa mãn  $F(0)=2$. Tính $F(-1)+2F(2)$.
	\loigiai{ 
		Ta  có $\displaystyle\int(2x-1)\mathrm{\,d}x=x^2-x+C_1; \displaystyle\int\left(3x^2-2\right)\mathrm{\,d}x=x^3-2x+C_2$.\\
		Suy ra $F(x)=\displaystyle\int f(x)\mathrm{\,d}x=\heva{&x^2-x+C_1\,\text{  khi } x \geq 1\\
		&x^3-2x+C_2\,\text{ khi  } x < 1}$ mà ta có $F(0)=2\Rightarrow C_2=2$.\\
		Mặt khác hàm số  $F$ là nguyên hàm của $f$ trên R nên $y=F(x)$ liên tục tại $x=1$.\\
		Suy  ra $\lim\limits_{x \to 1^{+}} F(x)=\lim\limits_{x \to 1^{-}} F(x) \Rightarrow C_1=1$.\\
		Khi  đó, ta có $F(x)=\heva{&x^2-x+1\text{ khi } x \geq 1\\
		&x^3-2x+2\text{  khi } x < 1}$ suy ra $\heva{&F(-1)=3\\
		 &F(2)=3.}$\\
		 Vậy $F(-1)+2F(2)=9$.
	 }
\end{bt}

% nguồn: Phan2_Lop_12_Dai_so_Chuong_4
\begin{bt}%[2D4H1-2]%[ma:976ae74d]
	Tìm hàm số $y=f(x)$  biết $f'(x)=3x^2+2x-m+1$, $f(2)=1$ và đồ thị của hàm số $y=f(x)$ cắt trục tung tại điểm có tung độ bằng $-5$.
	\loigiai{ 
		 Ta có $f(x)=\displaystyle\int\left(3x^2+2x-m+1\right)\mathrm{\,d}x=x^3+x^2+(1-m) x+C$.\\
		Theo đề bài, ta  có $\heva{&f(2)=1\\
		&f(0)=-5}  \Rightarrow\heva{&2(1-m)+C+12=1\\
		&C=-5}  \Rightarrow\heva{&m=4\\
		 &C=-5.}$\\
	Vây  $f(x)=x^3+x^2-3x-5$.}
\end{bt}
